\documentclass[runningheads,envcountsect]{lipics-v2021}

\usepackage{stmaryrd}
\usepackage{bbding}
%\Envelope
\newcommand{\cat}[1]{\ensuremath{\mathbf{#1}}}
\usepackage{amsthm,amsmath,amssymb,mathrsfs, dsfont}
\usepackage{hyperref, cleveref} %simboli matematici
\usepackage[all, cmtip]{xy}

\usepackage[utf8]{inputenc} % Input encoding - per caratteri particolari
%\usepackage[english]{babel} % Lingua principale inglese
\usepackage{graphicx} % Per includere immagini esterne
\usepackage[tickmarkheight=.5em,textwidth=\marginparwidth,textsize=small]{todonotes}
\usepackage{mathtools}
\usepackage{csquotes}


\usepackage{tikz-cd}
\usepackage{tikz}
\usetikzlibrary{decorations.markings}

% General Setting diagrams
%\usepackage{tikz-cd} %diagrammi
%\tikzcdset{row sep/normal=5em}
%\tikzcdset{column sep/normal=5em}
%\tikzcdset{every label/.append style = {font = \small}}

%\spnewtheorem*{notation}{Notation}{\bfseries}{}
%\spnewtheorem*{convention}{Convention}{\bfseries}{}



%funtori
\usepackage[usestackEOL]{stackengine}
\newcommand\functorop[1][l]{\csname#1functor\endcsname}
\newcommand\lfunctorop[3]{%
	\setbox0=\hbox{$#2$}%
	\kern\wd0%
	\ensurestackMath{\Centerstack[c]{#1\\ \mathllap{#2\;\,}\mathclap{\DownArrow}\\#3}}%
}		
\newcommand\rfunctorop[3]{%
	\setbox0=\hbox{$#2$}%
	\ensurestackMath{\Centerstack[c]{#1\\\mathclap{\UpArrow}\mathrlap{\,\;#2}\\#3}}%
	\kern\wd0%
}
\newcommand\functoropmapsto{\mathrel{\ensurestackMath{\Centerstack[c]{\longmapsto\\ \\\longmapsto}}}}
\setstackgap{L}{1.3\normalbaselineskip}
\newcommand\UpArrow{\rotatebox[origin=c]{90}{$\longrightarrow$\,}}
\newcommand\DownArrow{\rotatebox[origin=c]{-90}{$\longrightarrow$\,}}
\newcommand\functor[1][l]{\csname#1functor\endcsname}
\newcommand\lfunctor[3]{%
	\setbox0=\hbox{$#2$}%
	\kern\wd0%
	\ensurestackMath{\Centerstack[c]{#1\\ \mathllap{#2\;\,}\mathclap{\DownArrow}\\#3}}%
}
\newcommand\rfunctor[3]{%
	\setbox0=\hbox{$#2$}%
	\ensurestackMath{\Centerstack[c]{#1\\\mathclap{\DownArrow}\mathrlap{\,\;#2}\\#3}}%
	\kern\wd0%
}
\newcommand\functormapsto{\mathrel{\ensurestackMath{\Centerstack[c]{\longmapsto\\ \\\longmapsto}}}}
\setstackgap{L}{1.3\normalbaselineskip}

\newcommand{\lgh}{\mathsf{lg}}

\DeclareMathAlphabet{\mymathbb}{U}{BOONDOX-ds}{m}{n}

\newcommand{\Ob}{\mathcal{O}b}
\newcommand{\Hom}{\mathcal{H}om}
\newcommand{\Set}{\mathbf{Set}}
\newcommand{\Reg}{\mathcal{Reg}}
\newcommand{\Mono}{\mathcal{Mono}}
\newcommand{\initial}{\mymathbb{0}}
\newcommand{\terminal}{\mathds{1}}
\newcommand{\eg}[1]{\mathbf{EqGraph}_{\textbf {\textup{#1}}}}
\newcommand{\egg}[1]{\mathbf{EGG}_{\textbf {\textup{#1}}}}

\makeatletter
\def\@citecolor{blue}%
\def\@urlcolor{blue}%
\def\@linkcolor{blue}%
\def\UrlFont{\rmfamily}
\def\orcidID#1{\smash{\href{http://orcid.org/#1}{\protect\raisebox{-1.25pt}{\protect\includegraphics{orcid_color.eps}}}}}
\makeatother



\def\R{\mathsf{R}}
\def\B{\textbf {\textup{B}}}
\def\C{\textbf {\textup{C}}}
\def\D{\textbf {\textup{D}}}
\def\X{\textbf {\textup{X}}}
\def\Y{\textbf {\textup{Y}}}
\def\E{\textbf {\textup{E}}}
\def\T{\textbf {\textup{1}}}
\def\A{\textbf {\textup{A}}}
\def\M{\mathcal{M}}
%categorie varie
\newcommand{\catname}[1]{\textbf{\textup{#1}}}
\newcommand{\lab}{\catname{LHyp}}
\newcommand{\hyp}{\catname{Hyp}}
\newcommand{\hyps}{\catname{Hyp}_{\Sigma}}
\newcommand{\EqHyp}{\catname{EqHyp}} %equivalence hypergraphs
\newcommand{\gr}{\textbf{\textup{Graph}}}
\newcommand{\dgr}{\catname{SGraph}}
\newcommand{\dg}{\catname{DAG}}
\newcommand{\rt}{\mathsf{dcl_s}}
\newcommand{\rta}{\mathsf{dcl}}
\newcommand{\rtd}{\mathsf{dcl_{d}}}
\newcommand{\slice}[2]{(\catname{#1}\downarrow{#2})}
\newcommand{\tg}[0]{\catname{TG}_{\Sigma}}
\newcommand{\teg}[0]{\catname{TeGr}_{\Sigma}}
\newcommand{\sv}[0]{\mathsf{Sieves}}
\newcommand{\mono}[1]{\mathsf{Mon}(\catname{#1})}
\newcommand{\mo}[1]{{#1}_\mathsf{Mon}}
\newcommand{\pro}{\mathsf{prod}}
\newcommand{\spro}{\mathsf{ps}}
\newcommand{\prol}{\mathsf{lprod}}
\newcommand{\pred}[1]{{\downarrow}#1}
\newcommand{\colim}[0]{\mathrm{colim}}
\newcommand{\cod}{\mathsf{cod}}
\renewcommand{\sp}{\mathsf{sp}}
\renewcommand{\sup}{\mathsf{sup}}
\newcommand{\cow}[1]{\mathsf{cwd}({#1})}
\renewcommand{\inf}{\mathsf{inf}}
\newcommand{\dom}{\mathsf{dom}}
\newcommand{\dwnarrow}{\downarrow \hspace{-2pt}}
\newcommand{\Dwnarrow}{\Downarrow \hspace{-2pt}}

\newcommand{\ari}{\mathsf{ar}}
\newcommand{\abs}[1]{\lvert #1\rvert}


\newcommand{\upstr}[1] { {#1}^{\uparrow }}
\newcommand{\comma}[2]{#1\hspace{1pt} {\downarrow}#2}
\newcommand{\cma}[2]{\mathcal{#1}\hspace{1pt} {\downarrow}\hspace{1pt} \mathcal{#2}}


%sommatoria e prodotto per xy

\usepackage{relsize}
\newcommand{\Sum}{\mathlarger{\sum}}
\newcommand{\Prod}{\mathlarger{\prod}}

\bibliographystyle{abbrv}

%frecce
\newcommand{\mor}{\mathsf{Mor}}
\newcommand{\mon}{\mathsf{Mono}}
\newcommand{\reg}{\mathsf{Reg}}
\newcommand{\mto}{\rightarrowtail}
\newcommand{\eto}{\twoheadrightarrow}
\newcommand{\id}[1]{\mathsf{id}_{#1}}

%\tikzset{->-/.style={decoration={
%			markings,
%			mark=at position #1 with {\arrow{>}}},postaction={decorate}}}

%\titlerunning{????} %TODO optional, please use if title is longer than one line

%% Author with single affiliation.
%\author{Roberto Biondo \inst{1}{\small\Envelope}\and Davide Castelnovo\inst{1}\orcidID{0000-0002-5926-5615}
	%\and \\Fabio Gadducci\inst{1}\orcidID{0000-0003-0690-3051}
%}
%\institute{Dept.~of Computer Science, University of Pisa, Italy. 
	%\email{r.biondo@studenti.unipi.it},
	%\email{castelnovod@gmail.com},
	% \email{fabio.gadducci@unipi.it}
%}

\authorrunning{R.~Biondo, D.~Castelnovo, F.~Gadducci}
\Copyright{R.~Biondo, D.~Castelnovo, F.~Gadducci}

\ccsdesc[500]{Theory of Computation~Models of computation}
\ccsdesc[500]{Theory of Computation~Semantics and reasoning}

\keywords{Adhesive categories, EGGS} %TODO mandatory; please add comma-separated list of keywords

%Editor-only macros:: begin (do not touch as author)%%%%%%%%%%%%%%%%%%%%%%%%%%%%%%%%%%
\EventEditors{John Q. Open and Joan R. Access}
\EventNoEds{2}
\EventLongTitle{42nd Conference on Very Important Topics (CVIT 2016)}
\EventShortTitle{CVIT 2016}
\EventAcronym{CVIT}
\EventYear{2016}
\EventDate{December 24--27, 2016}
\EventLocation{Little Whinging, United Kingdom}
\EventLogo{}
\SeriesVolume{42}
\ArticleNo{23}


\title{On the adhesivity of EGGS}

\begin{document}
\maketitle

\begin{abstract}
\todo{a very nice abstract}
\end{abstract}


\section{Introduction}
\todo{A very nice introduction}
\section{$\mathcal{M}$-adhesive categories}


This  section recalls the basic theory of \emph{$\mathcal{M}$-adhesive categories} \cite{azzi2019essence,ehrig2012,ehrig2014adhesive,lack2005adhesive,heindel2009category}
and investigate some simple results for their relationship with kernel pairs.

%\begin{notation}
Given a category $\X$ we will not distinguish notationally between $\X$ and its class of objects, so
``$X\in \X$'' means that $X$ is an object of $\X$. We let $\mor(\X)$, $\mon(\X)$ and $\reg(\X)$ denote the class of all arrows, monos and regular monos of $\X$, respectively.  Given an object $X$, we will denote by $?_X$ the unique arrow from an initial object into $X$ and by $!_X$ that  unique arrow from $X$ into a terminal one.
%\end{notation}

\subsection{Background on $\mathcal{M}$-adhesivity}\label{subsec:ade}
The key property of $\mathcal{M}$-adhesive categories is the \emph{Van Kampen condition}~\cite{brown1997van,johnstone2007quasitoposes,lack2005adhesive}. 

In order to define it we need to introduce some terminology.  Let  $\X$ be a category. A subclass $\mathcal{A}$ of $\mor(\X)$ is called
\begin{itemize}
	\item		\emph{stable under pushouts (pullbacks)} if for every pushout (pullback) square as the one  below, if $m \in \mathcal{A}$ ($n\in \mathcal{A}$) then $n \in \mathcal{A}$ ($m \in \mathcal{A}$);
	\[\xymatrix{A \ar[r]^{f} \ar[d]_{m}& B \ar[d]^{n} \\ C \ar[r]_{g} & D}\]
		\item \emph{closed under composition} if $h, k\in \mathcal{A}$ implies $h\circ k\in \mathcal{A}$ whenever $h$ and $k$ are composable.
\end{itemize}

\begin{definition}
	Let $\mathcal{A}\subseteq \mor(\X)$ be a class of arrows in a category $\X$ and consider the following diagram.
	
		\[
	\xymatrix@C=10pt@R=6pt{&A'\ar[dd]|\hole_(.65){a}\ar[rr]^{f'} \ar[dl]_{m'} && B' \ar[dd]^{b} \ar[dl]_{n'} \\ C'  \ar[dd]_{c}\ar[rr]^(.7){g'} & & D' \ar[dd]_(.3){d}\\&A\ar[rr]|\hole^(.65){f} \ar[dl]^{m} && B \ar[dl]^{n} \\C \ar[rr]_{g} & & D }\]
	
	We say that the bottom square is an \emph{$\mathcal{A}$-Van Kampen square} if
		\begin{enumerate}
			\item it is a pushout square;
			\item 	whenever the cube above has pullbacks as back and left faces and the vertical arrows belong to $\mathcal{A}$, then its top face is a pushout 
			if and only if the front and right faces are pullbacks.
	\end{enumerate}


	Pushout squares that enjoy only the ``if'' half of item (2) above are called \emph{$\mathcal{A}$-stable}. A $\mor(\X)$-Van Kampen square is called  \emph{Van
		Kampen} and a $\mor(\X)$-stable square  \emph{stable}.
\end{definition}

We can now define $\mathcal{M}$-adhesive categories.

\begin{definition}
	Let $\X$ be a category and $\mathcal{M}$ a subclass of
	$\mon(\X)$  including  all isomorphisms, closed under composition,  and stable under pullbacks and pushouts.  The category  $\X$ is said to be \emph{$\mathcal{M}$-adhesive} if
	\begin{enumerate}
		\item it has \emph{$\mathcal{M}$-pullbacks}, i.e.~pullbacks along arrows of $\mathcal{M}$;
		\item it has \emph{$\mathcal{M}$-pushouts}, i.e.~pushouts along arrows of $\mathcal{M}$;
		\item  $\mathcal{M}$-pushouts are $\mathcal{M}$-Van Kampen squares.
	\end{enumerate}
	
	A category $\X$ is said to be \emph{strictly $\mathcal{M}$-adhesive}
	if $\mathcal{M}$-pushouts are Van Kampen. We write $m\colon X \mto Y$ to denote that an arrow $m\colon X\to Y$ belongs to $\mathcal{M}$.
\end{definition}



\begin{remark}
	\label{rem:salva}
	\emph{Adhesivity} and \emph{quasiadhesivity} 
	\cite{lack2005adhesive,garner2012axioms} coincide with strict
	$\mon(\X) $-adhesivity and strict $\reg(\X)$-adhesivity,
	respectively.
\end{remark}


$\mathcal{M}$-adhesivity is well-behaved with respect to  the construction of slice and functor categories \cite{mac2013categories}, as shown by the following theorems~\cite{ehrig2006fundamentals,lack2005adhesive}.

\begin{theorem}
	\label{thm:slice-functors}
	Let $\X$ be an (strict) $\mathcal{M}$-adhesive category, then the following hold true:
	\begin{enumerate}
		\item if $\Y$ is an (strict) $\mathcal{N}$-adhesive category $L\colon \Y\to \A$ a functor preserving $\mathcal{N}$-pushouts and $R\colon \X\to A$ one preserving pullbacks, then $\comma{L}{R}$ is (strictly) $\comma{\mathcal{N}}{\mathcal{M}}$-adhesive, where
		\[\comma{\mathcal{N}}{\mathcal{M}}:=\{(h,k) \in \mor(\comma{L}{R}) \mid h\in \mathcal{N}, k\in \mathcal{M}\}\]
		\item for every object $X$
		the categories $\X/X$  and $X/X$ are, respectively, (strictly) $\mathcal{M}/X$-adhesive and (strictly) $X/\mathcal{M}$-adhesive, where
		\[\mathcal{M}/X:=\{m\in \mor(\X/X) \mid m\in
		\mathcal{M}\} \hspace{5pt} X/\mathcal{M}:=\{m\in \mor(X/\X) \mid m\in \mathcal{M}\}\]
		\item for every small category $\Y$, the category $\X^\Y$ of
		functors $\Y\to \X$ is (strictly) $\mathcal{M}^{\Y}$-adhesive, where
		$\mathcal{M}^{\Y}:=\{\eta \in \mor(\X^\Y) \mid \eta_Y \in
		\mathcal{M} \text{ for every } Y\in \Y\}$;
		\item if $\Y$ is a full subcategory of $\X$ closed in it under pullbacks and $\mathcal{M}$-pushouts, then $\Y$ is (strictly) $\mathcal{N}$-adhesive for every class of arrows $\mathcal{N}$ of $\Y$ contained in $\mathcal{M}$ and  stable under pullbacks, pushout, containing all isomorphism and closed under composition and decomposition.
	\end{enumerate} 
\end{theorem}

We will briefly list some examples of $\mathcal{M}$-adhesive categories.

\begin{example}
	\label{ex:adhesive}
	$\cat{Set}$ is adhesive, and, more generally, every topos is
	adhesive~\cite{lack2006toposes}. By the closure properties above, every presheaf $[\cat{X},\cat{Set}]$ is adhesive, thus the category
	$\cat{Graph} = [ E \rightrightarrows V, \cat{Set}]$ is adhesive
	where $E \rightrightarrows {V}$ is the two objects category with two
	morphisms $s,t \colon{E} \to {V}$. Similarly, various
	categories of hypergraphs can be shown to be adhesive, such as term
	graphs and hierarchical graphs~\cite{CastelnovoGM24}. Note that the category $\cat{sGraphs}$ of simple graphs, 
	i.e.~graphs without parallel edges, is
	$\reg{(\cat{sGraphs})}$-adhesive~\cite{BehrHK23} but not
	quasiadhesive.
\end{example}


\begin{remark}\label{rem:deco}
	We can point out an important property of strict $\mathcal{M}$-adhesive categories with pullbacks and consider the $\mathcal{M}$-pushout square below.
	\[\xymatrix{A \ar[r]^{f} \ar@{>->}[d]_{m}& B \ar@{>->}[d]^{n}\\ C \ar[r]_{g} & D}\]
	
	Given an arrow $d\colon X\to D$, we can present the object $X$ has a pushout: indeed, consider the following cube, in which all the three vertical squares are pullbacks.
	\[
	\xymatrix@C=15pt@R=9pt{&V\ar[dd]|\hole_(.65){a}\ar[rr]^{v} \ar@{>.>}[dl]_{u} && Y \ar[dd]^{b} \ar@{>->}[dl]_{y} \\ Z  \ar[dd]_{c}\ar[rr]^(.7){z} & & X \ar[dd]_(.3){d}\\&A\ar[rr]|\hole^(.65){f} \ar@{>->}[dl]^{m} && B \ar@{>->}[dl]^{n} \\C \ar[rr]_{g} & & D }\]

Now, notice that, since the front square is a pullback, then the dotted arrow $u\colon V\to Z$ exists. Moreover, the usual composition and decomposition property of pullbacks \cite{mac2013categories} entails that the left face of the cube so obtained is a pullback too, proving that $u$ is in $\mathcal{M}$ and that the top square is a pushout.

	Clearly if the arrow $p\colon X\to D$ is in $\mathcal{M}$, we can omit the assumptions of strictness and the existence of all pullbacks.
\end{remark}


Finally, another  useful properties of $\mathcal{M}$-adhesive category are the following.
\begin{proposition}\label[proposition]{prop:regmono}
	If $\X$ is $\mathcal{M}$-adhesive then the following are true:
	\begin{enumerate}
		\item every $\mathcal{M}$-pushout square is also a pullback;
		\item every arrow in $\mathcal{M}$ is a regular mono.
	\end{enumerate}
\end{proposition}
\begin{proof}
	\begin{enumerate}
		\item  Consider the following cube in which the bottom face is an $\mathcal{M}$-pushout.
		\[\xymatrix@C=15pt@R=9pt{&A\ar[dd]|\hole_(.65){\id{A}}\ar[rr]^{g} \ar[dl]_{\id{A}} && B \ar[dd]^{\id{B}} \ar[dl]_(.6){\id{B}} \\ A  \ar@{>->}[dd]_{m}\ar[rr]^(.65){g} & & B \ar@{>->}[dd]_(.3){n}\\&A\ar[rr]|\hole^(.65){g} \ar@{>->}[dl]^{m} && B \ar@{>->}[dl]^{n} \\C \ar[rr]_{f} & & D}\]
		By construction the top face of the cube is a pushout and the back one a pullback. The left face is a pullback because $m$ is mono, thus the Van Kampen property yields that the front and the right faces are pullbacks too and the thesis follows.
		\item Let $m\colon X\mto Y$ be an arrow in $\mathcal{M}$, we can then take its pushout along itself, which, by the previous point, is also a pullback.
		\[\xymatrix{X \ar@{>->}[r]^{m} \ar@{>->}[d]_{m}& Y \ar@{>->}[d]^{h}\\ Y \ar@{>->}[r]^{k} & Z}\]
		It is now immediate to see that $m$ is the equalizer of $h$ and $k$.  
	\end{enumerate}
\end{proof}

\subsection{Some Results on Kernel Pairs and Regular Epis}

\begin{definition}[Kernel Pair]
    A \emph{kernel pair} for an arrow $f: A \to B$ is an object $K_f$ together with two arrows $\pi^1_f, \pi^2_f : K_f \to A$, denoted as $(K_f, \pi^1_f, \pi^2_f)$, such that the following square is a pullback.
    \[
        \begin{tikzcd}
            K_f \ar[r, "{\pi^1_f}"] \ar[d, "{\pi^2_f}" swap] & A \arrow[d, "f"] \\
            A \ar[r, "f"swap] & B
        \end{tikzcd}
    \]
\end{definition}

\begin{remark}
	If a category $\cat{C}$ has pullbacks then every arrow has a kernel pair.
\end{remark}

\begin{remark}
    Since a kernel pair is nothing more that a pullback, that is, a limit, by \Cref{rem:limits_are_unique_up_to_isomorphisms}, it make sense to refer to it as \emph{the} kernel pair for a morphism $f$.
\end{remark}

\begin{example}\label[example]{ex:kernel_pairs_in_Set}
    In $\Set$, a kernel pair for a function $f: A\to B$ is the set
    \[
        K_f=\{(x, y) \in A \times A \mid f(x) = f(y)\}
    \]
    together with the canonical projection on the first and the second component of the pairs.
\end{example}

\begin{proposition}\label[proposition]{prop:pairng_of_kernel_pairs_mono}
    Let $(K, p_1, p_2)$ be the kernel pair of $f: X \to Y$, and let $(X\times X, \pi_1, \pi_2)$ be the product of $X$ with itself. Then, the mediating arrow $\langle p_1, p_2\rangle : K \to X \times X$ is mono.
\end{proposition}

\begin{proof}
	Suppose $\langle p_1, p_2 \rangle \circ f = \langle p_1, p_2 \rangle \circ g$ for $f, g: Z \to K$. Then, we have
	\[
		\begin{split}
			\langle p_1, p_2 \rangle \circ f &=  \langle p_1, p_2 \rangle \circ g \\
			\pi_1 \circ \langle p_1, p_2 \rangle \circ f &=  \pi_1 \circ \langle p_1, p_2 \rangle \circ g \\
			p_1 \circ f &= p_1 \circ g
		\end{split}
		\qquad
		\begin{split}
			\langle p_1, p_2 \rangle \circ f &=  \langle p_1, p_2 \rangle \circ g \\
			\pi_2 \circ \langle p_1, p_2 \rangle \circ f &=  \pi_2 \circ \langle p_1, p_2 \rangle \circ g \\
			p_2 \circ f &= p_2 \circ g
		\end{split}
	\]
	Thus, from the universal property of the pullback, $f = g$.   
\end{proof}

\begin{proposition}\label[proposition]{prop:kermono}
	An arrow $m\colon M\to X$ is mono if and only if $(M, MM, \id{M})$ is a kernel pair for it.
\end{proposition}

\begin{proof}
    To prove the ``if'' part of the statement, let $f, g: A \to M$ be such that $m\circ f = m\circ g$, and consider the following situation.
    \[
        \begin{tikzcd}[row sep=26, column sep = 26]
        A \arrow[drr, bend left=30, "f"] \arrow[ddr, bend right=30,"g" swap] \arrow[dr, dashed, "u"] & & \\
        & M  \arrow[d, "{\id{M}}"] \arrow[r, "{\id{M}}"] & M \arrow[d, "m"] \\
        & M  \arrow[r, "m" swap] & X
        \end{tikzcd}
    \]
    For the universal property of pullbacks, we have that $$f  =  \id{M} \circ u =  g$$
    Hence, $m$ is mono.

    Conversely, if $m$ is mono, then, we have that
    \begin{align*}
        m \circ f = m \circ g   &\Rightarrow    f = g \\
                                &\Rightarrow    f \circ \id{M} = g\circ \id{M}
    \end{align*}
    Hence, $f$ is the unique arrow that makes the commutative square illustrated above a pushout.   
\end{proof}

\begin{remark}\label{rem:monos_in_presh_cats}
    From characterization of monos via pullbacks in \Cref{prop:kermono} and \Cref{lemma:limits_of_presheaves}, we have that a mono in a category of presheaves is a natural transformation of which each component is mono.
\end{remark}

\begin{corollary}\label[corollary]{cor:kermono}
	$(K_f, \pi_f^1, \pi_f^2)$ is a kernel pair for $f\colon X\to Y$ if and only if, for each mono $m\colon Y\to Z$, $(K_f, \pi_f^1, \pi_f^2)$ is a kernel pair also for $m\circ f$.
\end{corollary}
\begin{proof}
    It is enough to see that, by \Cref{lemma:pullback_lemma} and \Cref{prop:kermono} the outer boundary of the following square is a pullback.
        \[\begin{tikzcd}[row sep=13 pt, column sep=13 pt]
    	{K_f} && X && X \\
    	\\
    	X && Y && Y \\
    	\\
    	X && Y && Z
    	\arrow["{\pi_f^2}", from=1-1, to=1-3]
    	\arrow["{\pi_f^2}"', from=1-1, to=3-1]
    	\arrow["{id{X}}", from=1-3, to=1-5]
    	\arrow["f", from=1-3, to=3-3]
    	\arrow["f", from=1-5, to=3-5]
    	\arrow["f", from=3-1, to=3-3]
    	\arrow["{id_X}"', from=3-1, to=5-1]
    	\arrow["{id_Y}", from=3-3, to=3-5]
    	\arrow["{id_Y}", from=3-3, to=5-3]
    	\arrow["m", from=3-5, to=5-5]
    	\arrow["f"', from=5-1, to=5-3]
    	\arrow["m"', from=5-3, to=5-5]
    \end{tikzcd}\]
    The leftward part of the statement follows by definition of monomorphism an \Cref{lemma:pullback_lemma}.   
\end{proof}

\begin{lemma}\label[lemma]{lemma:kern_pairs_pres_pullbacks}
    Suppose the following situation, and that $f: X \to Y$ and $g: Z \to W$ have kernel pairs.
    \[
        \begin{tikzcd}
            X \ar[r, "h"] \ar[d, "f"swap] & Z \ar[d, "g"] \\
            Y \ar[r, "t"swap] & W
        \end{tikzcd}
    \]
    
    Then, there exists a unique arrow $k_h: K_f \to K_g$ making the squares below commute.
    \[
        \begin{tikzcd}[row sep = 25 pt, column sep= 25 pt]
            K_f \ar[r, dashed, "{k_h}"] \ar[d, "{\pi_f^1}"swap] & K_g \ar[d, "{\pi_g^1}"] \\
            X \ar[r, "h"swap] & Z 
        \end{tikzcd}
        \qquad
        \begin{tikzcd}[row sep = 25 pt, column sep= 25 pt]
            K_f \ar[r, dashed, "{k_h}"] \ar[d, "{\pi_f^2}"swap] & K_g \ar[d, "{\pi_g^2}"] \\
            X \ar[r, "h"swap] & Z 
        \end{tikzcd}
    \]

    Moreover, if the beginning square is a pullback, then also the preceding ones are so.
\end{lemma}

\begin{proof}
    Computing, we have
    \begin{align*}
        g \circ h \circ \pi_f^1 &=  t \circ f \circ \pi_f^1     \\
                                &=  t \circ f \circ \pi_f^2     \\
                                &=  g \circ h \circ \pi_f^2
    \end{align*}
    By the universal property of $K_g$ as the pullback of $g$ along itself, such $k_h$ exists and it is unique.

    To prove the second half of the thesis, let us consider the two rectangles below, which, by \Cref{lemma:pullback_lemma} are pullbacks.
    \[\begin{tikzcd}[row sep= 25 pt, column sep= 25 pt]
	{K_f} & X & Z \\
	X & Y & W
	\arrow["{\pi_f^1}", from=1-1, to=1-2]
	\arrow["{\pi_f^2}"', from=1-1, to=2-1]
	\arrow["h", from=1-2, to=1-3]
	\arrow["f", from=1-2, to=2-2]
	\arrow["g", from=1-3, to=2-3]
	\arrow["f"', from=2-1, to=2-2]
	\arrow["t"', from=2-2, to=2-3]
    \end{tikzcd}
    \qquad
    \begin{tikzcd}[row sep = 25 pt, column sep=25 pt]
	{K_f} & X & Z \\
	X & Y & W
	\arrow["{\pi_f^2}", from=1-1, to=1-2]
	\arrow["{\pi_f^1}"', from=1-1, to=2-1]
	\arrow["h", from=1-2, to=1-3]
	\arrow["f", from=1-2, to=2-2]
	\arrow["g", from=1-3, to=2-3]
	\arrow["f"', from=2-1, to=2-2]
	\arrow["t"', from=2-2, to=2-3]
    \end{tikzcd}
    \]
    But then the following ones are pullbacks too.
    \[\begin{tikzcd}[row sep= 25 pt, column sep= 25 pt]
	{K_f} & {K_g} & Z \\
	X & Y & W
	\arrow["{k_h}"', from=1-1, to=1-2]
	\arrow["{\pi_f^1}"', from=1-1, to=2-1]
	\arrow["{\pi_g^2}"', from=1-2, to=1-3]
	\arrow["{\pi_g^1}", from=1-2, to=2-2]
	\arrow["g", from=1-3, to=2-3]
	\arrow["h", from=2-1, to=2-2]
	\arrow["g", from=2-2, to=2-3]
        \arrow["{h\circ \pi_f^2}", from=1-1, to=1-3, bend left = 30]
        \arrow["{t \circ f}"', from=2-1, to=2-3, bend right = 30]
    \end{tikzcd}
    \qquad
    \begin{tikzcd}[row sep = 25 pt, column sep=25 pt]
	{K_f} & {K_g} & Z \\
	X & Y & W
	\arrow["{k_h}"', from=1-1, to=1-2]
	\arrow["{\pi_f^2}"', from=1-1, to=2-1]
	\arrow["{\pi_g^1}"', from=1-2, to=1-3]
	\arrow["{\pi_g^2}", from=1-2, to=2-2]
	\arrow["g", from=1-3, to=2-3]
	\arrow["h", from=2-1, to=2-2]
	\arrow["g", from=2-2, to=2-3]
        \arrow["{h\circ \pi_f^1}", from=1-1, to=1-3, bend left = 30]
        \arrow["{t \circ f}"', from=2-1, to=2-3, bend right = 30]
    \end{tikzcd}
    \]

    The thesis follows again by \Cref{lemma:pullback_lemma}.   
\end{proof}

\begin{proposition}\label[proposition]{prop:reg_epi_coeq}
    Let $e: X \to Y$ be a regular epimorphism in a category $\cat C$ with a kernel pair $(K, \pi_1, \pi_2)$. Then, $e$ is the coequalizer of $\pi_1$ and $\pi_2$.
\end{proposition}

\begin{proof}
    By hypothesis, there exists a pair $f, g: Z \to X$ of which $e$ is the coequalizer. Since $e \circ f = e \circ g$, we have
    \[
        \begin{tikzcd}[row sep= 20, column sep = 13]
            Z \ar[drr, "f", bend left=30] \ar[ddr, "g"swap, bend right=30] \ar[dr, dashed, "k"] & & \\
            & K \ar[r, "{\pi_1}"] \ar[d, "{\pi_2}" swap] & X \ar[d, "e"] \\
            & X \ar[r, "e"swap] & Y
        \end{tikzcd}
    \]
    thus there exists the unique $k: Z \to K$. Let now $h: Z \to V$ be an arrow such that $h \circ \pi_1 = h \circ \pi_2$, then
    \begin{align*}
        h \circ f &= h \circ \pi_1 \circ k \\
                  &= h \circ \pi_2 \circ k \\
                  &= h \circ g
    \end{align*}
    and thus there exists a unique $l: Y \to V$ such that $l \circ e = h$.   
\end{proof}

\begin{corollary}\label[corollary]{cor:reg_epi_components_reg_epi_nat_trans}
    Let $\cat C$ be a category with pullbacks and $\phi : D \dot\to D'$ be a natural transformation between two functors $D, D': \cat{I \to C}$. If $\phi_i$ is a regular epi for every $i$, then $\phi$ is a regular epi.
\end{corollary}

\begin{proof}
    Let $(K_i, \pi_i^1, \pi_i^2)$ be the kernel pair of $\phi_i$ for each $i$. Given an arrow $\alpha: i \to j$ of $\cat I$, we have
    \begin{align*}
        \phi_j \circ D(\alpha) \circ \pi_i^1 &= D'(\alpha) \circ \phi_i \circ \pi_i^1 \\
                                             &= D'(\alpha) \circ \phi_i \circ \pi_i^2 \\
                                             &= \phi_j \circ D(\alpha) \circ \pi_i^2
    \end{align*}
    Thus, the outer boundary of the diagram below commutes, yielding the arrow $K(\alpha)$
    \[
        \begin{tikzcd}[row sep=26, column sep=26]
            {K_i} \ar[r, "{\pi_i^1}"] \ar[d, "{\pi_i^2}"swap] \ar[dr, dashed, "{K(\alpha)}"] & D(i) \ar[dr, "{D(\alpha)}"] \\
            D(i) \ar[dr, "{D(\alpha)}"swap] & K_j \ar[r, "{\pi_j^1}"] \ar[d, "{\pi_j^2}"swap] & D(j) \ar[d, "{\phi_j}"] \\
            & D(j) \ar[r, "{\phi_j}"swap] & D'(j)
        \end{tikzcd}
    \]

    In this way, we get a functor $K: \cat{I \to C}$, 
    which maps each $i$ onto $K_i$ and each arrow $\alpha$ onto $K(\alpha)$. 
	We have in fact $K(id_i) : K_i \to K_i$ is the arrow such that
        \[
            \begin{split}
                D(id_i) \circ \pi_i^1 &= \pi_i^1 \circ K(id_i) \\
                \pi_i^1 &= \pi_i^1 \circ K(id_i)
            \end{split}
                \qquad
            \begin{split}
                D(id_i) \circ \pi_i^2 &= \pi_i^2 \circ K(id_i) \\
                \pi_i^2 &= \pi_i^2 \circ K(id_i)
            \end{split}
        \]
	Thus, for the universal property of pullbacks, $K(id_i) = id_{K_i}$.

        Suppose now $\alpha : i \to j$ and $\beta: j \to k$. Computing, we have
        \[
            \begin{split}
                \pi_k^1 \circ K(\beta \circ \alpha) &= D(\beta \circ \alpha) \circ \pi_i^1 \\
                                                    &= D(\beta) \circ D(\alpha) \circ \pi_i^1 \\
                                                    &= D(\beta) \circ \pi_j^1 \circ K(\alpha) \\
                                                    &= \pi_k^1 \circ K(\beta) \circ K(\alpha)
            \end{split} \qquad
            \begin{split}
                \pi_k^2 \circ K(\beta \circ \alpha) &= D(\beta \circ \alpha) \circ \pi_i^2 \\
                                                    &= D(\beta) \circ D(\alpha) \circ \pi_i^2 \\
                                                    &= D(\beta) \circ \pi_j^2 \circ K(\alpha) \\
                                                    &= \pi_k^2 \circ K(\beta) \circ K(\alpha)
            \end{split}
        \]
        Again, for universal property of pullbacks, necessarily we have $K(\beta \circ \alpha) = K(\beta) \circ K(\alpha)$, proving functoriality of $K$.
   
    
     Hence, we have two natural transformations $\pi^1, \pi^2 : E \dot\to D$. By \Cref{prop:reg_epi_coeq}, every component $\phi_i$ is the coequalizer of $\pi_i^1, \pi_i^2: E \to D$, and so $\phi$ is the coequalizer of $\pi^1$ and $\pi^2$.   
\end{proof}

\begin{lemma}\label[lemma]{lemma:nat_trans_reg_epi_canonical_arrow_reg_epi}
    Let $D, D': \cat{I \to C}$ be two diagrams, and let $((c_i)_{i \in \cat I}, C)$ and $((c_i')_{i\in \cat I}, C')$ be, respectively, the colimit of $D$ and $D'$. If $\cat C$ has all colimits, for diagrams of shape $\cat I$ and $\phi: D \dot\to D'$ is a natural transformation in which all components are regular epimorphisms, then, the arrow induced by $\phi$ from $C$ to $C'$ (\Cref{prop:nat_tran_induces_unique_arrow_between_colimits}) is a regular epimorphism too.
\end{lemma}

\begin{proof}
    By \Cref{cor:reg_epi_components_reg_epi_nat_trans}, we know that $\phi: D \dot\to D'$ is a regular epimorphism, so that there is a functor $E: \cat{I \to C}$ and $\eta, \theta: E \dot\to D$ such that $\phi$ is the coequalizer of $\eta$ and $\theta$. Let now $((p_i)_{i \in \cat I}, P)$ be the colimit of $E$, by \Cref{prop:nat_tran_induces_unique_arrow_between_colimits}, we have $a, b: P \to C$ fitting in the diagram below.
    \[
        \begin{tikzcd}[row sep = 26, column sep = 26]
            E(i) \ar[r, "{p_i}"] \ar[d, "{\eta_i}"swap] & P \ar[d, dashed, "a"] \\
            D(i) \ar[r, "{c_i}"swap] & C
        \end{tikzcd}
        \qquad
        \begin{tikzcd}[row sep = 26, column sep = 26]
            E(i) \ar[r, "{p_i}"] \ar[d, "{\theta_i}"swap] & P \ar[d, dashed, "b"] \\
            D(i) \ar[r, "{c_i}"swap] & C
        \end{tikzcd}
    \]

    We want to show that $c$ coequalizes $\eta$ and $\theta$. Let thus $t: C \to T$ be an arrow such that $t \circ a = t \circ b$. Then, for every $i$, we have
    \begin{align*}
        t \circ c_i \circ \eta_i &= t \circ a \circ p_i \\
                                 &= t \circ b \circ p_i \\
                                 &= t \circ c_i \circ \theta_i
    \end{align*}    

    Thus, there is $t_i: D(i) \to T$ such that $t\circ c_i = t_i \circ \phi_i$. It is now easy to see that $((t_i)_{i \in \cat I}, T)$ is a cocone of $D'$: suppose $\alpha: i \to j$ be an arrow of $\cat I$, obtaining
    \begin{align*}
        t_i \circ \phi_i &= t \circ c_i \\
                         &= t \circ c_j \circ D(\alpha) \\
                         &= t_j \circ \phi_j \circ D(\alpha) \\
                         &= t_j \circ D'(\alpha) \circ \phi_i
    \end{align*}
    By the hypothesis that $\phi_i$ is regular epi for each $i$, therefore epi (by the dual of \Cref{prop:eq_are_mono}), we can conclude $t_i = t_j \circ D'(\alpha)$.
    
    Hence, we have an arrow $k: C' \to T$ such that $k \circ c_i' = t_i$. But then
    \begin{align*}
        c \circ c \circ c_i &= k \circ c_i' \circ \phi_i \\
                            &= t_i \circ \phi \\
                            &= t \circ c_i
    \end{align*}
    Showing that $k \circ c = t$.

    For the uniqueness, let $k': C' \to T$ be another arrow such that $k' \circ c = t$. Then we have
    \begin{align*}
        k' \circ c_i' \circ \phi_i &= k' \circ c \circ c_i \\
                                   &= t \circ c_i \\
                                   &= t_i \circ \phi_i
    \end{align*}    
    Since $\phi_i$ is a regular epimorphism, we have $k' \circ c_i' = t_i$, and, because $k \circ c_i' = t_i$ by construction, we can conclude that $k'=k$ since $((c_i')_{i \in \cat I}, C')$ is a colimit.   
\end{proof}


\iffalse 
\subsection{Factorization systems}

\begin{definition}[Factorization System]\label{def:fact_sys}
	Let $\X$ be a category, and let $\E$, $\M$ be two classes of arrows.
	Then, $(\E, \M)$ is a \emph{factorization system} if:
	\begin{enumerate}
		\item $\E$ and $\M$ are closed under composition with isomorphisms: if $f: X \to Y$ belongs to $\E$ (to $\M$), and $h: Y\to Z$ is an isomorphism, then $h \circ f$ belongs to $\E$ (to $\M$);
		\item every arrow $f: X \to Y$ of $\X$ admits an \emph{$(\E, \M)$-factorization}, that is, there are arrows $e \in \E$ and $m \in \M$ such that $f = m\circ e$;
		\item every $e \in \E$ has the \emph{left lifting property} with respect to every $m \in \M$: for every commutative square as the one below, with $e \in \E$ and $m \in \M$,
			there exists a unique $k: Y \to Z$ whitch fits in it.
			\[\xymatrix{
					X\ar[r]^{g} \ar[d]_e & Z \ar[d]^m \\ Y \ar@{.>}[ur]^k \ar[r]_f & V
			}\]
	\end{enumerate}
	A factorization system $(\E, \M)$ is \emph{proper} if every $e \in \E$ is epi and every $m \in \M$ is mono, while it is \emph{stable} is $\E$ is stable under pullbacks.
\end{definition}

\begin{remark}\label{rem:dual_fact_sys}
	It is immediate to notice that, given a factorization system $(\E, \M)$ on a category $\X$, then $(\M, \E)$ is a factorization system on $\X^{op}$, which is proper if $(\E, \M)$ is so, and it is stable if
	$\M$ is stable under pushouts in $\X$.
\end{remark}

\begin{remark}
	If $(\E, \M)$ is a factorization system for a category $\X$, then $\E$ is the class of all the arrows with the left lifting property with respect to every $m\in \M$.
	Indeed, if $f:X \to Y$ is an arrow with such a property with respect to every element of $\M$, it is possible to factorize it as $f = m \circ e$, with $m: M \to Y$ in $\M$ and 
	$e: X \to M$ in $\E$, obtaining the above situation by definition.
	\[\xymatrix{X \ar[r]^e\ar[d]_f & M \ar[d]^m \\ Y \ar[r]_{\id{Y}} \ar@{.>}[ur]^k & Y}\]
	From this, it is possible to prove that $k$ and $m$ are mutually inverses, hence $f = m\circ e$ belongs to $\E$ due to closure under isomorphisms.
\end{remark}

\begin{proposition}
	Let $(\E, \M)$ be a factorization system on $\X$. Then $\E$ contains every isomorphism, and it is closed under composition.
\end{proposition}

By \Cref{rem:dual_fact_sys}, we have also that $\M$ is closed under composition, and contains every isomorphism.

\begin{lemma}
	Let $(\E, \M)$ be a factorization system on $\X$. If $e:Y \to E$, $e': Y\to E'$ and $m: E \to X$, $m':E'\to X$ are arrows, respectively, in $\E$ and in $\M$ such that $e'\circ m' = e \circ m$,
	then there exists a unique isomorphism $\phi: E \to E'$ such that the diagram below commutes.
	\[\xymatrix{Y\ar[r]^{e'}\ar[d]_{e}&E'\ar[d]^{m'}\\E\ar[r]_{m}\ar@{.>}[ur]^{\phi}\ar[r]_{m}&X}\]
\end{lemma}

\begin{proof}
	Suppose the following situation.
	\[\xymatrix{X \ar[r]^e\ar[d]_{e'}&E\ar[d]^{m}&&X\ar[r]^{e'}\ar[d]_{e}&{E'}\ar[d]^{m'}\\{E'}\ar[r]_{m'}\ar@{.>}[ur]^{\phi}&Y&&E\ar[r]_{m}\ar@{.>}[ur]^{\psi}&Y}\]
	From the commutativity of the diagrams below, we obtain at once $\phi\circ\psi=\id{E'}$ and $\psi\circ\phi=\id{E}$.
	\[\xymatrix{
			X\ar[rr]^{e}\ar[dd]_{e}\ar[dr]_{e'}&&E\ar[r]^{m}\ar[dl]^{\psi}&Y\ar[dd]^{\id{Y}} &&& X\ar[rr]^{e'}\ar[dd]_{e'}\ar[dr]_{e}&&{E'}\ar[r]^{m}\ar[dl]^{\phi}&Y\ar[dd]^{\id{Y}}\\
							     & {E'}\ar[dl]_{\phi}\ar[drr]^{m'}& & &&&  & {E'}\ar[dl]_{\psi}\ar[drr]^{m'}& & \\
			E\ar[rrr]_{m} &&& Y &&& {E'} \ar[rrr]_{m'} &&& Y
	}\]
\end{proof}

The previous lemma allows us to deduce other useful properties of factorization systems.

\begin{proposition}
	Given a factorization system $(\E, \M)$ on a category $\X$, it holds that:
	\begin{enumerate}
		\item an arrow $f:X\to Y$ that factors as $m \circ e$ belongs to $\M$ if and only if $e$ is an isomorphism;
		\item $f\in \E$ and $f\in\M$ if and only if $f$ is an isomorphism;
		\item $\M$ is stable under pullbacks;
		\item if $(\E, \M)$ is proper, then $g\circ f \in \M$ implies $f \in \M$.
	\end{enumerate}
\end{proposition}

\todo{proof?}
\fi 




\section{Hypergraphical structures}\label{sec:hyper}

In this section we briegly recall the notion \emph{$\X$-hypergraph}. In order to do it is necessary to have a monad $(-)^\star$ playing a role analogous to the usual \emph{Kleene star}, also known as \emph{list monad}, $(-)^\star\colon \Set\to \Set$ sending a set to the free monoid on it \cite{sakarovitch2009elements,wadler1995monads}. We will recall some of its proprieties.

\begin{proposition}
	The following facts are true:
	\begin{enumerate}
		\item for every set $X$ and $n\in \mathbb{N}$ there are arrows $v_{n}\colon X^n\to X^\star $ such that $(X^\star, \{v_{n}\}_{n\in \mathbb{N}})$ is a coproduct;
		\item for every arrow $f\colon X\to Y$, $f^\star$ is the coproduct of the family $\{f^n\}_{n\in \mathbb{N}}$;
		\item $(-)^\star$ preserves all \emph{connected limits} \cite{carboni1995connected}, in particular it preserves all pullbacks and equalizers.
	\end{enumerate}
\end{proposition}

\begin{remark}\label{rem:mono}
	Preservation of pullbacks implies that $(-)^\star$ sends monos to monos.
\end{remark}

\begin{remark}\label[remark]{rem:length}Notice that, $1^\star$ can be canonically identified with $\mathbb{N}$, thus for every set $X$, the arrow $!_\X\colon X\to 1$ induces a \emph{length function} $!^\star_X\colon X^\star \to \mathbb{N}$ which simply sends a word length.
\end{remark}

\subsection{The category of hypergraphs}

We will start this section with the definition of hypergraphs and we will see how to label them with an algebraic signature.  


\begin{definition}An \emph{hypergraph} is a 4-uple $\mathcal{G}:=(E_\mathcal{G}, V_\mathcal{G}, s_\mathcal{G}, t_\mathcal{G})$ made by two sets $E_\mathcal{G}$ and $V_\mathcal{G}$, called respectively the set of \emph{hyperedges} and \emph{nodes}, plus a pair of \emph{source} and \emph{target} arrows  $s_\mathcal{G}, t_\mathcal{G}\colon E_\mathcal{G}\rightrightarrows V_\mathcal{G}^\star$. A \emph{hypergraph morphism} $(E_\mathcal{G}, V_\mathcal{G}, s_\mathcal{G}, t_\mathcal{G})\to (E_\mathcal{H}, V_\mathcal{H}, s_\mathcal{H}, t_\mathcal{H})$ is a pair $(h,k)$ of functions $h\colon E_\mathcal{G}\to E_\mathcal{H}$, $k\colon V_\mathcal{G}\to V_\mathcal{H}$ such that the following diagrams commute.
	
	\[\xymatrix{ E_{\mathcal{G}} \ar[d]_{h} \ar[r]^{s_{\mathcal{G}}}& V^\star_{\mathcal{G}}  \ar[d]^{k^\star}& E_{\mathcal{G}} \ar[r]^{t_{\mathcal{G}}} \ar[d]_{h} & V^\star_{\mathcal{G}} \ar[d]^{k^\star}  \\ E_{\mathcal{G}} \ar[r]_{s_{\mathcal{H}}} & V^\star_{\mathcal{H}} & E_{\mathcal{G}} \ar[r]_{t_{\mathcal{H}}} & V^\star_{\mathcal{H}} }\]
	We define $\hyp$ to be the resulting category.
\end{definition}

Let $\pro^\star$ be the functor sending $X$ to $X^\star\times X^\star$, then we can present $\hyp$ as a comma category.
\begin{proposition}\label[proposition]{prop:com}
	$\hyp$ is isomorphic to $\comma{\id{\Set}}{\pro^\star}$
\end{proposition}

\begin{corollary}\label[corollary]{prop:hypadh}
	 $\hyp$ is an adhesive category.
\end{corollary}
\begin{proof}
	By hypothesis $(-)^\star$ preserves pullbacks, while $\pro$ is continuous by definition, thus the thesis follows from  \Cref{thm:slice-functors} and \Cref{prop:com}.   
\end{proof}

Another useful corollary of \Cref{prop:com} is the following one.

\begin{corollary}\label[corollary]{cor:monhyper}
A morphism $(h,k)$ is a mono in $\hyp$ if and only if both its components are injective functions.
\end{corollary}

\Cref{prop:left,prop:com} allows us to deduce immediately the following.

\begin{proposition}\label[proposition]{cor:left}  The forgetful functor $U_{\hyp}$ which sends an hypergraph $\mathcal{G}$ to its object of nodes has a left adjoint $\Delta_{\hyp}$.
\end{proposition}

\begin{example}Since the initial object of $\catname{Set}$ is the empty set,  $\Delta_{\Set}(X)$ is the hypergraph which has $X$ as set of nodes and $\emptyset$ as set of hyperedges and $?_X$ as both source and target function.
\end{example}

\todo{forse possiamo anche cancellare tutti gli esempi sotto}

We can represent hypergraphs graphically. We will use dots to denote nodes and squares to denote hyperedges, the name of a node or of an hyperedge will be put near the corresponding dot or square. Sources and targets are represented by lines between dots and squares: the lines from the sources of an hyperedge will have an arrowhead in the middle pointing towards the hyperedge, while the lines to the targets will have arrowheads pointing to the target nodes.  We will decorate the arrow corresponding to the $i^{th}$ letter of a target or a source with a label $i$.


\begin{example}Take $V_{\mathcal{G}}$ to be be $\{v_1, v_2, v_3, v_4, v_5\}$ and $E_{\mathcal{G}}$ to be $\{h_1, h_2, h_3\}$. Sources and targets are given by:
	\[\begin{matrix}
		s_{\mathcal{G}}(h_1)\colon 2\to V_{\mathcal{G}}  & \begin{matrix}
			0 \mapsto v_1\\
			1\mapsto v_2
		\end{matrix} && s_{\mathcal{G}}(h_2)\colon 2\to V_{\mathcal{G}} & \begin{matrix}
			0 \mapsto v_3\\
			1\mapsto v_4 
		\end{matrix} && s_{\mathcal{G}}(h_3)\colon 1\to V_{\mathcal{G}} & 
		0 \mapsto v_5\\
		t_{\mathcal{G}}(h_1)\colon 2\to V_{\mathcal{G}} & \begin{matrix}
			0 \mapsto v_3\\
			1\mapsto v_4
		\end{matrix} && t_{\mathcal{G}}(h_2)\colon 2\to V_{\mathcal{G}} & 0\mapsto v_5 && t_{\mathcal{G}}(h_3)\colon 0\to V_{\mathcal{G}} &  t_{\mathcal{G}}(h_3)=?_{ V_{\mathcal{G}}} 
	\end{matrix}\]
	
	We can draw the resulting $\mathcal{G}$ as follows:
	\begin{center}\begin{tikzpicture}
			\node[circle,fill=black,inner sep=0pt,minimum size=6pt,label=above:{$v_1$}] (A) at (0,0) {};
			\node[circle,fill=black,inner sep=0pt,minimum size=6pt,label=above:{$v_2$}] (B) at (0,-1.5) {};
			\node[circle,fill=black,inner sep=0pt,minimum size=6pt,label=above:{$v_3$}] (C) at (3,0) {};
			\node[circle,fill=black,inner sep=0pt,minimum size=6pt,label=above:{$v_4$}] (D) at (3,-1.5) {};
			\node[circle,fill=black,inner sep=0pt,minimum size=6pt,label=above:{$v_5$}] (E) at (6,-0.75) {};
			\draw[rounded corners] (1.25, -1) rectangle (1.75, -0.5) {};
			\draw[->-=.5](4.75,-0.75)--(E)node[pos=0.5, above,font=\fontsize{7}{0}\selectfont]{$1$};
			\draw[->-=.5](E)--(7,-0.75)node[pos=0.5, above,font=\fontsize{7}{0}\selectfont]{$1$};
			\draw[rounded corners] (4.25, -1) rectangle (4.75, -0.5) {};
			\node at (4.5, -0.3){$h_2$};
			\node at (1.5, -0.3){$h_1$};
			\node at (7.25, -0.3){$h_3$};
			\draw[rounded corners] (7, -1) rectangle (7.5, -0.5) {};
			\draw(A)[->-=.5]..controls(0.5,0)and(1.2,-0.2)..(1.25,-0.6)node[pos=0.5, above,font=\fontsize{7}{0}\selectfont]{$1$};
			\draw(B)[->-=.5]..controls(0.5,-1.5)and(1.2,-1.3)..(1.25,-0.9)node[pos=0.5, below,font=\fontsize{7}{0}\selectfont]{$2$};
			
			\draw(C)[->-=.5]..controls(3.5,0)and(4.2,-0.25)..(4.25,-0.6)node[pos=0.5, above,font=\fontsize{7}{0}\selectfont]{$1$};
			\draw(D)[->-=.5]..controls(3.5,-1.5)and(4.2,-1.3)..(4.25,-0.9)node[pos=0.5, below,font=\fontsize{7}{0}\selectfont]{$2$};
			
			\draw[->-=.5](1.75,-0.9)..controls(1.8,-1.3)and(2.5,-1.5)..(D)node[pos=0.5, below,font=\fontsize{7}{0}\selectfont]{$2$};
			\draw[->-=.5] (1.75,-0.6)..controls(1.8,-0.25)and(2.5,0)..(C) node[pos=0.5, above,font=\fontsize{7}{0}\selectfont]{$1$};
		\end{tikzpicture}
	\end{center}
\end{example}
\begin{example}\label[example]{exa_2} Let $V_{\mathcal{G}}$ be as in the previous example and $E_{\mathcal{G}}=\{h_1, h_2, h_3\}$.	Then we define
	\[\begin{matrix}
		s_{\mathcal{G}}(h_1)\colon 0\to V_{\mathcal{G}}  & s_{\mathcal{G}}(h_1)=?_{V_\mathcal{G}} && s_{\mathcal{G}}(h_2)\colon 2\to V_{\mathcal{G}} & \begin{matrix}
			0 \mapsto v_1\\
			1\mapsto v_2
		\end{matrix}&& s_{\mathcal{G}}(h_3)\colon 2\to V_{\mathcal{G}} & \begin{matrix} 
			0 \mapsto v_1\\
			1\mapsto v_4	
		\end{matrix}\\
		t_{\mathcal{G}}(h_1)\colon 1\to V_{\mathcal{G}} & 
		0 \mapsto v_1 && 		t_{\mathcal{G}}(h_2)\colon 1\to V_{\mathcal{G}} & 0\mapsto v_3 &&   t_{\mathcal{G}}(h_3)\colon 1\to V_{\mathcal{G}} & 1\mapsto v_5
	\end{matrix}\]
	
	Now we can depict $\mathcal{G}$ as
	\begin{center}\begin{tikzpicture}
			\node[circle,fill=black,inner sep=0pt,minimum size=6pt,label=above:{$v_1$}] (A) at (0,0) {};
			\node[circle,fill=black,inner sep=0pt,minimum size=6pt,label=above:{$v_2$}] (B) at (0,-1.5) {};
			\node[circle,fill=black,inner sep=0pt,minimum size=6pt,label=above:{$v_3$}] (C) at (3,-0.75) {};
			\node[circle,fill=black,inner sep=0pt,minimum size=6pt,label=above:{$v_4$}] (D) at (3,-2.25) {};
			\node[circle,fill=black,inner sep=0pt,minimum size=6pt,label=right:{$v_5$}] (E) at (6,-1.5) {};
			\draw[->-=.5] (1.75,-0.75)--(C)node[pos=0.5, above,font=\fontsize{7}{0}\selectfont]{$1$};
			\draw[rounded corners] (1.25, -1) rectangle (1.75, -0.5) {};
			\draw[->-=.5] (4.75,-1.5)--(E)node[pos=0.5, above,font=\fontsize{7}{0}\selectfont]{$1$};
			\draw[->-=.5] (-1.5,0)--(A)node[pos=0.5, above,font=\fontsize{7}{0}\selectfont]{$1$};
			\draw[rounded corners] (4.25, -1.75) rectangle (4.75, -1.25) {};
			\node at (4.5, -1.05){$h_3$};
			\node at (1.5, -0.3){$h_2$};
			\node at (-1.75, 0.45){$h_1$};
			\draw[rounded corners] (-2, -0.25) rectangle (-1.5, 0.25) {};
			\draw[->-=.5] (A)..controls(0.5,0)and(1.2,-0.2)..(1.25,-0.6)node[pos=0.5, above,font=\fontsize{7}{0}\selectfont]{$1$};
			\draw[->-=.5] (B)..controls(0.5,-1.5)and(1.2,-1.3)..(1.25,-0.9)node[pos=0.5, below,font=\fontsize{7}{0}\selectfont]{$2$};
			
			\draw[->-=.5] (C)..controls(3.5,-0.75)and(4.2,-0.95)..(4.25,-1.35)node[pos=0.5, above,font=\fontsize{7}{0}\selectfont]{$1$};
			\draw[->-=.5] (D)..controls(3.5,-2.25)and(4.2,-2.05)..(4.25,-1.65)node[pos=0.5, below,font=\fontsize{7}{0}\selectfont]{$2$};
		\end{tikzpicture}
	\end{center}
\end{example}

\subsubsection{$\hyp$ as a category of functors}

Following \cite{bonchi2022string}, we can present $\hyp$ as a category of functor over a suitable category.

\begin{definition}Let $\catname{H}$ be the category in which:
	\begin{itemize}
		\item the set of objects is given by $ (\mathbb{N}\times \mathbb{N}) \cup \{\bullet\}$
		\item arrows are given by the identities $\id{k,l}$ and $\id{\bullet}$ and exactly $k+l$ arrows $f_i\colon (k,l)\rightarrow \bullet$, where $i$ ranges from $0$ to $k+l-1$;
		\item composition is defined simply putting, for every $f_i\colon (k,l)\rightarrow \bullet$:
		\begin{equation*}
			f_i=f_i\circ \id{k,l} \qquad f_i = \id{\bullet}\circ f_i 
		\end{equation*}
	\end{itemize}
\end{definition}

The idea now is that for every functor $F\colon \catname{H}\to \catname{Set}$ we can define
\[E_F:=\sum_{k,l\in \mathbb{N}}F(k,l)\]

Now, for every $k$, $l$, $i$ and $j$ in $\mathbb{N}$ with $i< k$ and $j< l$ we can define $s^F_{k,l}\colon F(k,l)\to F(\bullet)^k$,  $t^F_{k,l}\colon F(k,l)\to F(\bullet)^l$ as the unique ones fitting in the diagrams below, where the vertical arrows are the projections.
\[\xymatrix{F(k,l)  \ar@{.>}[r]^{s^F_{k,l}} \ar[dr]_{F(f_i)}& F(\bullet)^{k} \ar[d]^{\pi^F_{k,i}} & F(k,l) \ar@{.>}[r]^{t^F_{k,l}} \ar[dr]_{F(f_{k+j})} & F(\bullet)^{l} \ar[d]^{\pi^F_{l,j}} \\ & F(\bullet) && F(\bullet)}\]

In turn, these arrows allow us to consider
$s_F, t_F\colon E_F\rightrightarrows F(\bullet)^{\star}$ as the unique arrows fitting in the diagrams below, where the vertical arrows are coprojections.
\[\xymatrix{F(k,l) \ar[d]_{a^F_{k,l}}  \ar[r]^{s^{F}_{k,l}}& F(\bullet)^{k} \ar[d]^{b^F_{k}} & F(k,l) \ar[d]_{a^F_{k,l}}  \ar[r]^{t^{F}_{k,l}}& F(\bullet)^{l} \ar[d]^{b^F_{l}}\\ E_F \ar@{.>}[r]_-{s_F}& F(\bullet)^\star & E_F \ar@{.>}[r]_-{t_F}& F(\bullet)^\star}\]

Let $\mathcal{G}_F$ be the resulting hypergraph. One can now show that  sending $F$ to $\mathcal{G}_F$ can be extended to an equivalence $\mathcal{G}_{-}\colon \Set^{\catname{H}}\to \hyp$ (see \cite{castelnovo2023thesis, castelnovo2022newcriterionmathcalmmathcalnadhesivity} for details).


\begin{proposition}
	$\hyp$ is equivalent to the category $\Set^{\catname{H}}$.
\end{proposition}

\subsubsection{Labelling hypergraph with an algebraic signature}

Our interest intohypergraph stems fromm their use as a graphical representation of algebraic terms. We thus need a way to label hyperedges with symbols taken from a given algebraic signature.

\begin{definition}
An \emph{algebraic signature} $\Sigma$ is a pair $(O_\Sigma, \ari_\Sigma)$ given by a \emph{set of operations} $O_\Sigma$ and an \emph{arity function} $\ari_\Sigma\colon O_\Sigma \to \mathbb{N}$. 

We define the \emph{hypergraph $\mathcal{G}_\Sigma$ associated with $\Sigma$} taking $O_\Sigma$ as set of hyperedges, $1$ as setof nodes, so that $1^\star$ is $\mathbb{N}$, $\ari_\sigma$ as the source function and $\delta_1$ as target function, where $\delta_1$ picks the element $1$. The category $\hyps$ of \emph{algebraically labelled hypergraphs} is the slice category $\hyp/\mathcal{G}^\Sigma$.
\end{definition}


\begin{example}\label[example]{exa_3} Let $\Sigma=(O_\Sigma, \ari_\Sigma)$ be an algebraic signature in $\Set$. This simply amount to a set of \emph{operations} with an associated natural number, called \emph{arity}. 	For instance let $\Sigma_G$ be the signature of groups, then $\mathcal{G}^{\Sigma_G}$ can be depicted as:
	\begin{center}
		\begin{tikzpicture}
			\node[circle,fill=black,inner sep=0pt,minimum size=6pt,label=above:{$v$}] (V) at (0,0) {};
			\node(E)at(-2, 0.4){$e$};
			\node(M)at(0, 2.15){$\cdot$};
			\node(I)at(2, 0.5){$(-)^{-1}$};
			\draw[->-=.5](-1.75,0)--(V)node[pos=0.5, above,font=\fontsize{7}{0}\selectfont]{$1$};
			\draw[->-=.5](V)..controls(-0.5,0.5)and(-0.8,1)..(-0.25,1.6)node[pos=0.5, right,font=\fontsize{7}{0}\selectfont]{$2$};
			\draw[->-=.5](V)..controls(-1,0.6)and(-1,1.1)..(-0.25,1.9)node[pos=0.5, left,font=\fontsize{7}{0}\selectfont]{$1$};
			\draw[->-=.5](0.25,1.75)..controls(0.8,0.8)and(0.5,0.5)..(V)node[pos=0.5, right,font=\fontsize{7}{0}\selectfont]{$1$};
			\draw[->-=.5](V)--(1.75,0)node[pos=0.5, above,font=\fontsize{7}{0}\selectfont]{$1$};
			\draw[->-=.5](2.25,0)..controls(3.5,0)and(2.5,-2)..(V)node[pos=0.5, below,font=\fontsize{7}{0}\selectfont]{$2$};
			\draw[rounded corners] (-2.25, -0.25) rectangle (-1.75, 0.25) {};
			\draw[rounded corners] (-0.25, 1.5) rectangle (0.25, 2) {};
			\draw[rounded corners] (2.25, -0.25) rectangle (1.75, 0.25) {};
		\end{tikzpicture}
	\end{center}
\end{example}

\Cref{cor:mono} and \Cref{thm:slice-functors} give us immediately an adhesivity result for $\hyp_{\Sigma}$ and a characterization of monomorphisms in it.

\begin{proposition}\label[proposition]{prop:mono} For every algebraic signature $\Sigma$, the following are true:
	\begin{enumerate}
		\item a morphism $(h,k)$ between two object of $\hyp_{\Sigma}$ is a mono if and only if $h$ and $k$ are injective functions;
		\item $\hyp_{\Sigma}$ is an adhesive category. 
	\end{enumerate}
\end{proposition}


\begin{remark}\label[remark]{rem:label}	
	Let $\mathcal{H}=(E, V, s, t)$ be an hypergraph, by definition we know that $U_{\hyp}(\mathcal{G}^{\Sigma})$ is the terminal object $1$, so an arrow $\mathcal{H}\rightarrow \mathcal{G}^{\Sigma}$, is determined by a morphism $h\colon E_\mathcal{H}\to O_\Sigma$  making the two squares below commute (cfr.~ \Cref{rem:length}).
	\[\xymatrix{E_{\mathcal{H}} \ar[r]^{h} \ar[d]_{s_{\mathcal{H}}}& O_\Sigma \ar[d]^{\ari_\Sigma}& E_{\mathcal{H}} \ar[r]^{h} \ar[d]_{t_{\mathcal{H}}} & O_\Sigma \ar[d]^{\delta_1}\\ V^\star_{\mathcal{{H}}} \ar[r]_{\lgh_{V_{\mathcal{H}}}}& \mathbb{N} & V^\star_{\mathcal{{H}}} \ar[r]_{\lgh_{V_{\mathcal{H}}}} & \mathbb{N}}\]
	
	Let $v_n\colon V^n_\mathcal{H}\to  V^\star_{\mathcal{H}}$ be a coprojection.  The second diagram in particular above entails that $t_{\mathcal{H}}$ factors through the inclusion $v_1\colon V_{\mathcal{H}}\to V^{\star}$ of words of length $1$, i.e. that $t_{\mathcal{H}}=v_1\circ \tau_{\mathcal{H}}$ for some $\tau_{\mathcal{H}}\colon E_{\mathcal{H}}\to \tau_{\mathcal{H}}$.
\end{remark}

$\hyp_{\Sigma}$, has a forgetful functor $U_{\Sigma}\colon \hyp_{\Sigma}\to \X$ which sends $(h,k)\colon \mathcal{H}\to \mathcal{G}^{\Sigma}$ to $U_{\X}(\mathcal{H}$). Now, $U_{\X}(\mathcal{G}^{\Sigma})=1$ thus, for every object $X$, there is only one arrow $X\to U_{\X}(\mathcal{G}^{\Sigma})$. Define $\Delta_{\Sigma}(X)\colon \Delta_{\X}(X)\to \mathcal{G}^{\Sigma}$ as the transpose of this arrow. Explicitly, $\Delta_{\X}(X)=(0, X, ?_{X^\star}, ?_{X^\star})$ and $\Delta_\Sigma(X)$ is simply $(?_{O_\Sigma}, !_{X})$.

\begin{proposition} $U_\Sigma$
	has a left adjoint $\Delta_\Sigma$.
\end{proposition}
\begin{proof}Let $(h, !_{V_\mathcal{H}})\colon \mathcal{H}\to \mathcal{G}^{\Sigma}$ be an object of $\hyp_{\Sigma}$, and suppose that there exists $f\colon X\to U_{\Sigma}(\mathcal{H})$. Since, $U_{\Sigma}(\mathcal{H})=U_{\X}(\mathcal{H})$ and the identity is the unit of $\Delta_\hyp \dashv U_{\hyp}$, we get a morphism $(?_{E_{\mathcal{H}}},f)\colon \Delta_{\X}(X)\to \mathcal{H}$ of $\hyp$. But then we have
	\[
	(h, !_{V_{\mathcal{{H}}}})\circ (?_{E_{\mathcal{H}}}, f)=(h\circ ?_{E_{\mathcal{H}}}, !_{V_{\mathcal{{H}}}}\circ f)=(?_{0_{\Sigma}}, !_{X})=\Delta_{\hyp}(X)\]
and the thesis follow.   
\end{proof}

\todo{anche questo forse val la pena toglierlo}
We will extend our graphical notation of hypergraphs to labeled ones putting the label of an hyperedge $h$ inside its corresponding square.
\begin{example}\label[example]{lab_1}
	The simplest example is given by the identity $\id{\mathcal{G}^\Sigma}\colon \mathcal{G}^\Sigma\rightarrow \mathcal{G}^{\Sigma}$. If $\Sigma$ is the signature of groups $\Sigma_G$we get 
	\begin{center}
		\begin{tikzpicture}
			\node[circle,fill=black,inner sep=0pt,minimum size=6pt,label=above:{$\heartsuit$}] (V) at (0,0) {};
			\node(E)at(-2, 0.4){$e$};
			\node(M)at(0, 2.15){$\cdot$};
			\node(I)at(2, 0.75){$(-)^{-1}$};
			
			\node(E')at(-2, 0){$e$};
			\node(M')at(0, 1.75){$\cdot$};
			\node(I')at(2, 0){$(-)^{-1}$};
			\draw[->-=.5](-1.75,0)--(V)node[pos=0.5, above,font=\fontsize{7}{0}\selectfont]{$1$};
			\draw[->-=.5](V)..controls(-0.5,0.5)and(-0.8,1)..(-0.25,1.6)node[pos=0.5, right,font=\fontsize{7}{0}\selectfont]{$2$};
			\draw[->-=.5](V)..controls(-1,0.6)and(-1,1.1)..(-0.25,1.9)node[pos=0.5, left,font=\fontsize{7}{0}\selectfont]{$1$};
			\draw[->-=.5](0.25,1.75)..controls(0.8,0.8)and(0.5,0.5)..(V)node[pos=0.5, right,font=\fontsize{7}{0}\selectfont]{$1$};
			\draw[->-=.5](V)--(1.5,0)node[pos=0.5, above,font=\fontsize{7}{0}\selectfont]{$1$};
			\draw[->-=.5](2.5,0)..controls(4,0)and(2.5,-2)..(V)node[pos=0.5, below,font=\fontsize{7}{0}\selectfont]{$2$};
			\draw[rounded corners] (-2.25, -0.25) rectangle (-1.75, 0.25) {};
			\draw[rounded corners] (-0.25, 1.5) rectangle (0.25, 2) {};
			\draw[rounded corners] (2.5, -0.5) rectangle (1.5, 0.5) {};
		\end{tikzpicture}
	\end{center}
\end{example}

\begin{example}\label[example]{lab_2}
	Take again $\Sigma_G$ the signature of groups, then the hypergraph $\mathcal{G}$ of \Cref{exa_2} can be labeled defining
	\begin{align*}
		e=f(h_1) \quad \cdot=f(h_2)\quad \cdot=f(h_3)
	\end{align*}
	In this case we get the following picture
	
	\begin{center}\begin{tikzpicture}
			\node[circle,fill=black,inner sep=0pt,minimum size=6pt,label=above:{$v_1$}] (A) at (0,0) {};
			\node[circle,fill=black,inner sep=0pt,minimum size=6pt,label=above:{$v_2$}] (B) at (0,-1.5) {};
			\node[circle,fill=black,inner sep=0pt,minimum size=6pt,label=above:{$v_3$}] (C) at (3,-0.75) {};
			\node[circle,fill=black,inner sep=0pt,minimum size=6pt,label=above:{$v_4$}] (D) at (3,-2.25) {};
			\node[circle,fill=black,inner sep=0pt,minimum size=6pt,label=right:{$v_5$}] (E) at (6,-1.5) {};
			\draw[->-=.5] (1.75,-0.75)--(C)node[pos=0.5, above,font=\fontsize{7}{0}\selectfont]{$1$};
			\draw[rounded corners] (1.25, -1) rectangle (1.75, -0.5) {};
			\draw[->-=.5] (4.75,-1.5)--(E)node[pos=0.5, above,font=\fontsize{7}{0}\selectfont]{$1$};
			\draw[->-=.5] (-1.5,0)--(A)node[pos=0.5, above,font=\fontsize{7}{0}\selectfont]{$1$};
			\draw[rounded corners] (4.25, -1.75) rectangle (4.75, -1.25) {};
			\node at (4.5, -1.05){$h_3$};
			\node at (1.5, -0.3){$h_2$};
			\node at (-1.75, 0.45){$h_1$};
			\draw[rounded corners] (-2, -0.25) rectangle (-1.5, 0.25) {};
			\draw[->-=.5] (A)..controls(0.5,0)and(1.2,-0.2)..(1.25,-0.6)node[pos=0.5, above,font=\fontsize{7}{0}\selectfont]{$1$};
			\draw[->-=.5] (B)..controls(0.5,-1.5)and(1.2,-1.3)..(1.25,-0.9)node[pos=0.5, below,font=\fontsize{7}{0}\selectfont]{$2$};
			
			\draw[->-=.5] (C)..controls(3.5,-0.75)and(4.2,-0.95)..(4.25,-1.35)node[pos=0.5, above,font=\fontsize{7}{0}\selectfont]{$1$};
			\draw[->-=.5] (D)..controls(3.5,-2.25)and(4.2,-2.05)..(4.25,-1.65)node[pos=0.5, below,font=\fontsize{7}{0}\selectfont]{$2$};
			
			\node at (-1.75,0) {$e$};
			\node at (1.5,-0.75) {$\cdot$};
			\node at (4.5,-1.5) {$\cdot$};
		\end{tikzpicture}
	\end{center}
\end{example}


\subsection{Term Graphs}
\todo{a very nice introduction}

Let us start using labelled hypergraphs to define term graphs.

\begin{definition}\label[definition]{def:tg}\index{term graph}
	Given an algebraic signature $\Sigma$, we say that a labelled hypergraph $(l, !_{V_\mathcal{G}})\colon \mathcal{G}\to \mathcal{G}^{\Sigma}$ is a \emph{term graph} if $t_\mathcal{G}$ is mono. We define $\tg$ to be the full subcategory of $\hyp_{\Sigma}$ and denote by $I_\Sigma$ the inclusion. Restricting $U_\Sigma\colon \hyp_{\Sigma}\to \catname{Set}$ we get a forgetful functor $U_{\tg}\colon \tg\to \catname{Set}$.
\end{definition}


\begin{remark}\label[remark]{rem:mono2}By \Cref{rem:label}, we know that if $\mathcal{G}$ is a term graph then $t_{\mathcal{G}}=v_1\circ \tau_{\mathcal{G}}$, where $v_1$ is the coprojection of $V_{\mathcal{G}}$ into $V^\star_{\mathcal{G}}$.  Notice, moreover that since $t_{\mathcal{G}}$ is mono then $\tau_{\mathcal{G}}$ is mono too.
\end{remark}

We are now going back to examine the properties of $\tg$, with the purpose of studying its adhesivity properties.

\begin{proposition}\label[proposition]{term:left}The forgetful functor $U_{\tg}$ has a left adjoint $\Delta_{\tg}$.
\end{proposition}
\begin{proof}
	This follows noticing that, for every object $X$, $\Delta_{\Sigma}(X)$ is a term graph.   
\end{proof}



We canlist some categorical properties of $\tg$


\begin{proposition}
For every algebraic signatur $\Sigma$, the following are true:
\begin{enumerate}
	\item if  $(i,j)\colon \mathcal{H}\to \mathcal{G}$ is a mono between  $(l, !_{V_\mathcal{G}})\colon \mathcal{G}\to \mathcal{G}^{\Sigma}$ and $(l', !_{V_\mathcal{H}})\colon \mathcal{H}\to \mathcal{G}^{\Sigma}$ in $\hyp_\Sigma$ and the latter is a term graph, then also the former is in $\tg$
	\item $\tg$ has equalizers, binary products and pullbacks and they are created by $I_\Sigma$.
\end{enumerate}
\end{proposition}

\begin{remark}
	$\tg$ in general does not have terminal objects. Consider an algebraic signature in $\Set$. Since $U_{\tg}$ preserves limits, if a terminal object exists it must have the singleton as set of nodes, therefore the set of hyperedges must be empty or a singleton $\{h\}$. Now take as signature the one given by two operations $a$ and $b$, both of arity $0$; we have three term graphs with only one node $v$: $\Delta_{\tg{\Sigma}}(\{v\})$, $(l_a, !_{V_{\mathcal{G}}})\colon \mathcal{G}_a\to \mathcal{G}^{\Sigma}$ and $(l_b, !_{V_{\mathcal{G}}})\colon \mathcal{G}_b\to \mathcal{G}^{\Sigma}$.
	\begin{center}\begin{tikzpicture}
			
			\node[circle,fill=black,inner sep=0pt,minimum size=6pt,label=left:{$v$}] (V) at (5,0) {};
			\node[circle,fill=black,inner sep=0pt,minimum size=6pt,label=left:{$v$}] (U) at (3,0) {};
			\node at(5,1.25){$a$};	
			\node at(5,1.7){$h$};	
			
			\draw[->-=.5](5,1)--(V)node[pos=0.5, right,font=\fontsize{7}{0}\selectfont]{$1$};
			
			\draw[rounded corners] (4.75, 1) rectangle (5.25, 1.5) {};
			
			
			\node[circle,fill=black,inner sep=0pt,minimum size=6pt,label=left:{$v$}] (V) at (7,0) {};
			
			\node at(7,1.25){$b$};	
			\node at(7,1.7){$h$};	
			
			\draw[->-=.5](7,1)--(V)node[pos=0.5, right,font=\fontsize{7}{0}\selectfont]{$1$};
			
			\draw[rounded corners] (6.75, 1) rectangle (7.25, 1.5) {};
			
		\end{tikzpicture}
	\end{center}
	There are no morphisms in $\tg$ between the last two and from the last two to the first one, therefore none of them can be terminal.
\end{remark}

\begin{remark}
	$\tg$ is not an adhesive category. In particular it does not have pushouts along all monomorphisms. For instance, if we take the three term graphs of the previous remark, then have two arrows
	$(?_{\{h\}}, \id{\{v\}})\colon \Delta_{\tg}(\{v\})\to (l_a, !_{V_{\mathcal{G}_a}})$ and $(?_{\{h\}}, \id{\{v\}})\colon \Delta_{\tg}(\{v\})\to (l_b, !_{V_{\mathcal{G}_a}})$ which cannot be completed to a square. Indeed if $(q, !_{V_\mathcal{H}})\colon \mathcal{H}\to \mathcal{G}^\Sigma$ is another term graph with $(g_E, g_V)\colon (l_a, !_{V_{\mathcal{G}}})\to (q, !_{V_\mathcal{H}})$ and $(k_E, k_V)\colon (l_a, !_{V_{\mathcal{G}}})\to (q, !_{V_\mathcal{H}})$  such that 
	\[(g_E, g_V)\circ (?_{\{h\}}, \id{\{v\}}) = (k_E, k_V)\circ (?_{\{h\}}, \id{\{v\}})\]
	then $g_V=k_V$ and
	\[t_{\mathcal{H}}(g_E(h))=g^\star_V(t_{\mathcal{G}}(h))=g_V^\star(\delta_v)=k^\star_V(\delta_V)=k^\star_V(t_{\mathcal{G}}(h))=t_{\mathcal{H}}(k_E(h))\]
	so that we also have $g_E=k_E$, but then
	\[
	a=l_a(h)=q(g_E(h))=q(k_E(h))=l_b(h)=b\]
\end{remark}


\begin{definition}
	Let $(l, !_{V_{\mathcal{G}}})\colon \mathcal{G}\to \mathcal{G}^{\Sigma}$  be a term graph. A \emph{input node} is an element of $V_{\mathcal{G}}$ not in the image of $\tau_{\mathcal{H}}$.  A morphism $(f,g)$ between
	Let $(l, !_{V_{\mathcal{G}}})\colon \mathcal{G}\to \mathcal{G}^{\Sigma}$ and $(l, !_{V_{\mathcal{H}}})\colon \mathcal{H}\to \mathcal{G}^{\Sigma}$ in $\tg$, is said to \emph{preserve input nodes} if $g$ sends input nodes to input nodes.
\end{definition}

\todo{Se questo remark sotto non serve nel pezzo sulle equivalenze possiamo toglierlo}
\begin{remark}\label{prop:image}
	Suppose that $(f,g)\colon ((l, !_{V_{\mathcal{G}}}))\to (l', !_{V_{\mathcal{H}}})$ preserves input nodes. Then  if $\tau_{\mathcal{H}}(h)=g(v)$ for some $v\in V_{\mathcal{G}}$ then $h$ belongs to the image of $f$. Indeed, by hypothesis $v$ must be in the image of $\tau_{\mathcal{G}}$ and so there exists $k$ such that $\tau_{\mathcal{G}}(k)=v$. But then $\tau_{\mathcal{H}}(f(k))=g(v)$ and we can conclude that $f(k)=h$.
\end{remark}

Preservation of inputs, characterizes regular monos in $\tg$.

\begin{proposition}\label[proposition]{lem:reg} The following are equivalent for a a mono $(i,j)$ between two term graphs  $(l, !_{V_{\mathcal{G}}})\colon \mathcal{G}\to \mathcal{G}^{\Sigma}$ and  $(l', !_{V_{\mathcal{H}}})\colon \mathcal{H}\to \mathcal{G}^{\Sigma}$:
	\begin{enumerate}
		\item $(i, j)$ is a regular mono;
		\item $(i,j)$ preserves input nodes.
	\end{enumerate}	
\end{proposition}

This characterization, in turn, provides us with the following result \cite{castelnovo2022newcriterionmathcalmmathcalnadhesivity,castelnovo2023thesis}. 

\begin{lemma}\label[lemma]{prop:push} Consider three term graphs $(l_0, !_{V_\mathcal{G}})\colon \mathcal{G}\to \mathcal{G}^{\Sigma}$, $(l_1, !_{V_\mathcal{H}})\colon \mathcal{H}\to \mathcal{G}^{\Sigma}$ and $(l_2, !_{V_\mathcal{K}})\colon \mathcal{K}\to \mathcal{G}^{\Sigma}$. Given $(f_1, g_1)\colon (l_0, !_{V_\mathcal{G}})\to (l_1, !_{V_\mathcal{H}})$, $(f_2, g_2)\colon (l_0, !_{V_\mathcal{G}})\to (l_2, !_{V_\mathcal{K}})$, if $(f_1, g_1)$ is a regular mono, then its pushout along $(f_2, g_2)$,  then their pushout $(p, !_{V_{\mathcal{P}}})\colon \mathcal{P}\to \mathcal{G}^{\Sigma}$ in $\hyp_{\Sigma}$ is a term graph too.
\end{lemma}


\Cref{thm:slice-functors}, \Cref{lem:reg} and \Cref{prop:push} allow us to recover the following result, previously proved by direct computation in \cite[Thm.~4.2]{corradini2005term}.
\begin{corollary}
	The category $\tg$ is quasiadhesive.
\end{corollary}


\subsection{Hierarchical hypergraphs}

\todo{da fare? lasciare per una versione estesa?}


\commentato{
\section{Equipping hypergraphical structures with equivalences}


\subsection{Hypergraphs with equivalences}

\subsection{Term graphs with equivalences}

\todo{Il materiale va distribuito nelle varie sottosezioni}

\begin{definition}
	A \emph{hypergraph with equivalence} $\mathcal{G} = (E_\mathcal{G}, V_{\mathcal{G}}, C_\mathcal{G}, s_\mathcal{G}, t_\mathcal{G}, q_\mathcal{G})$ is a 6-tuple such that $\mathcal{G} = (E_\mathcal{G}, V_{\mathcal{G}}, s_\mathcal{G}, t_\mathcal{G})$ is a hypergraph, $C_\mathcal{G}$ is an object and $q_{\mathcal{G}}: V_{\mathcal{G}}\to C_{\mathcal{G}}$ is a regular epimorphism called \emph{quotient map}. 
	
	A morphism $h:\mathcal{G\to H}$ is a triple $(h_E, h_V, h_C)$ such that the following diagrams commute.
	\[\xymatrix{
		{E_\mathcal{G}}\ar[r]^{s_\mathcal{G}}\ar[d]_{h_E} & {V_{\mathcal{G}}^\star}\ar[d]^{h_V^\star} & {E_\mathcal{G}}\ar[r]^{t_\mathcal{G}}\ar[d]_{h_E} & {V_{\mathcal{G}^\star}}\ar[d]^{h_V^\star} & {V_\mathcal{G}}\ar[r]^{q_\mathcal{G}}\ar[d]_{h_V} & {C_{\mathcal{G}}} \ar[d]^{h_C} \\
		{E_\mathcal{H}}\ar[r]_{s_\mathcal{H}} & {V_{\mathcal{H}}^\star}	& {E_\mathcal{H}}\ar[r]_{t_\mathcal{H}} & {V_{\mathcal{H}^\star}}& {V_\mathcal{H}}\ar[r]_{q_\mathcal{H}} & {C_\mathcal{H}}
	}\]
	The category of hypergraphs with equivalences and their morphisms is denoted $\EqHyp$.

\end{definition}

\begin{remark}\label{rem:eqhyp_morphs}
	Morphisms of hypergraphs with equivalences are uniquely determined by the first two components. That is, if $h_1 = (h_E, h_V, f)$ and $h_2 = (h_E, h_V, g)$ are two morphisms $\mathcal{G \to H}$, then we have
	\[\xymatrix{
			{V_\mathcal{G}} \ar[r]^{h_V}\ar[d]_{q_\mathcal{G}} & V_\mathcal{H} \ar[d]^{q_\mathcal{H}} & V_\mathcal{G}\ar[l]_{h_V}\ar[d]^{q_\mathcal{G}}\\
			C_{\mathcal{G}}\ar[r]_{f} & C_{\mathcal{H}} & C_{\mathcal{G}}\ar[l]^{g}
	}\]
	Hence $f \circ q_\mathcal{G} = q_\mathcal{H}\circ h_V  =g\circ q_\mathcal{G}$, and since $q_\mathcal{G}$ is epi, we obtain $f = g$.
\end{remark}

$\EqHyp$ has a forgetful functor $U_{\EqHyp}:\EqHyp \to \Set$, which sends each $\mathcal{G} = (E_\mathcal{G}, V_{\mathcal{G}}, C_\mathcal{G}, s_\mathcal{G}, t_\mathcal{G}, q_\mathcal{G})$ into $V_\mathcal{G}$, and each $h = (h_E, h_V, h_C)$ onto $h_V$. 

\begin{proposition}
	$U_\EqHyp$ has a left adjoint $\Delta_{\EqHyp}: \Set \to \EqHyp$.
\end{proposition}

\begin{proof}
	For each set $X$, define $\Delta_\EqHyp(X):= (\emptyset, X, \{\bullet\}, ?_X, ?_X, !_X)$. Consider now $h: \Delta_{\EqHyp}(X) \to \mathcal{H}$.
	\[\xymatrix@C=2.3cm{
			\Delta_{\EqHyp}(X) \ar@{.>}[d]_{\Delta_\EqHyp(f)} \ar[dr]^{h} & \\
			\Delta_\EqHyp(U_{\EqHyp}(\mathcal{H})) \ar[r]_{\epsilon_{\mathcal{H}}} & \mathcal{H}
	}\]
	Where $\Delta_\EqHyp(U_{\EqHyp}(\mathcal{H})) = (\emptyset, V_\mathcal{H}, \{\bullet\}, ?_{V_\mathcal{H}}, ?_{V_\mathcal{H}}, !_{V_\mathcal{H}})$ and $\epsilon_{\mathcal{H}} = (?_{E_\mathcal{H}}, \id{V_\mathcal{H}}, g)$.
	Note that, since $\Delta_{\EqHyp}(X)$ has the empty set as object of edges, $h_E = ?_{E_\mathcal{H}}$, then, the unique arrow that fits in the diagram is $\Delta_{\EqHyp}(f) = (?_{E_\mathcal{H}}, h_V, \id{\{\bullet\}})$.

\end{proof}

We now define another functor $T: \EqHyp \to \hyp$, which ``forgets'' the quotient part, mapping each hypergraph with equivalence $\mathcal{G} = (E_\mathcal{G}, V_{\mathcal{G}}, C_\mathcal{G}, s_\mathcal{G}, t_\mathcal{G}, q_\mathcal{G})$ onto $T(\mathcal{G})=(E_{\mathcal{G}}, V_{\mathcal{G}}, s_\mathcal{G}, t_{\mathcal{G}})$. Then, we have the following result.

\begin{proposition}
	$T$ has a left adjoint $L: \hyp \to \EqHyp$.
\end{proposition}

\begin{proof}
	Let $\mathcal{G}$ be a hypergraph, and define $L(\mathcal{G}) := (E_\mathcal{G}, V_{\mathcal{G}}, \{\bullet\}, s_\mathcal{G}, t_\mathcal{G}, !_{V_\mathcal{G}})$. Let now $h: L(\mathcal{G})\to \mathcal{H}$ be a morphism in $\EqHyp$, and consider the following situation.
	\[\xymatrix@C=2.3cm{
		L(\mathcal{G}) \ar@{.>}[d]_{L(f)} \ar[dr]^{h}\\ L(T(\mathcal{H})) \ar[r]_{\epsilon_{\mathcal{H}}} & \mathcal{H}
	}\]
	for $L(T(\mathcal{H}))=(E_\mathcal{H}, V_{\mathcal{H}}, \{\bullet\}, s_\mathcal{H}, t_\mathcal{H}, !_{V_\mathcal{H}})$. Then $\epsilon_\mathcal{H} = (\id{E_\mathcal{H}}, \id{V_\mathcal{H}}, h_C)$ (by \Cref{rem:eqhyp_morphs}, the last component is uniquely determined by the first two), and $L(f)$ must be $(h_E, h_V, \id{\{\bullet\}})$.
\end{proof}

\begin{remark}
	$T$ is faithful. Indeed, consider two morphisms $h = (h_E, h_V, h_C)$ and $k = (k_E, k_V, k_C)$, and suppose $T(h) = T(k)$, that is, $(h_E, h_V) = (k_E, k_V)$.
	By \Cref{rem:eqhyp_morphs}, we can conclude also $h_C = k_C$, and hence the faithfulness of $T$.
\end{remark}
}

\section{Hypergraphs and term graphs with equivalences}
\label{hypereq}

\begin{definition}
	A \emph{hypergraph with equivalence} $\mathcal{G} = (E_\mathcal{G}, V_{\mathcal{G}}, C_\mathcal{G}, s_\mathcal{G}, t_\mathcal{G}, q_\mathcal{G})$ is a 6-tuple such that $\mathcal{G} = (E_\mathcal{G}, V_{\mathcal{G}}, s_\mathcal{G}, t_\mathcal{G})$ is a hypergraph, $C_\mathcal{G}$ is a set and $q_{\mathcal{G}}: V_{\mathcal{G}}\to C_{\mathcal{G}}$ is a regular epi called \emph{quotient map}. 
	%
	A morphism $h:\mathcal{G\to H}$ is a triple $(h_E, h_V, h_C)$ such that the following diagrams commute
	\[\xymatrix{
		{E_\mathcal{G}}\ar[r]^{s_\mathcal{G}}\ar[d]_{h_E} & {V_{\mathcal{G}}^\star}\ar[d]^{h_V^\star} & {E_\mathcal{G}}\ar[r]^{t_\mathcal{G}}\ar[d]_{h_E} & {V_{\mathcal{G}^\star}}\ar[d]^{h_V^\star} & {V_\mathcal{G}}\ar[r]^{q_\mathcal{G}}\ar[d]_{h_V} & {C_{\mathcal{G}}} \ar[d]^{h_C} \\
		{E_\mathcal{H}}\ar[r]_{s_\mathcal{H}} & {V_{\mathcal{H}}^\star}	& {E_\mathcal{H}}\ar[r]_{t_\mathcal{H}} & {V_{\mathcal{H}^\star}}& {V_\mathcal{H}}\ar[r]_{q_\mathcal{H}} & {C_\mathcal{H}}
	}\]
	The category of hypergraphs with equivalences and their morphisms is denoted $\EqHyp$.

\end{definition}

\begin{remark}\label{rem:eqhyp_morphs}
	Morphisms of hypergraphs with equivalences are uniquely determined by the first two components. That is, if $h_1 = (h_E, h_V, f)$ and $h_2 = (h_E, h_V, g)$ are two morphisms $\mathcal{G \to H}$, then we have
	\[\xymatrix{
			{V_\mathcal{G}} \ar[r]^{h_V}\ar[d]_{q_\mathcal{G}} & V_\mathcal{H} \ar[d]^{q_\mathcal{H}} & V_\mathcal{G}\ar[l]_{h_V}\ar[d]^{q_\mathcal{G}}\\
			C_{\mathcal{G}}\ar[r]_{f} & C_{\mathcal{H}} & C_{\mathcal{G}}\ar[l]^{g}
	}\]
	Hence $f \circ q_\mathcal{G} = q_\mathcal{H}\circ h_V =g\circ q_\mathcal{G}$, and since $q_\mathcal{G}$ is epi, we obtain $f = g$.
\end{remark}

$\EqHyp$ has a forgetful functor $U_{\EqHyp}:\EqHyp \to \Set$, which sends each $\mathcal{G} = (E_\mathcal{G}, V_{\mathcal{G}}, C_\mathcal{G}, s_\mathcal{G}, t_\mathcal{G}, q_\mathcal{G})$ into $V_\mathcal{G}$, and each $h = (h_E, h_V, h_C)$ onto $h_V$. 

\begin{proposition}
	$U_\EqHyp$ has a left adjoint $\Delta_{\EqHyp}: \Set \to \EqHyp$.
\end{proposition}

\begin{proof}
	For each set $X$, define $\Delta_\EqHyp(X):= (\emptyset, X, \{\bullet\}, ?_X, ?_X, !_X)$. Consider now $h: \Delta_{\EqHyp}(X) \to \mathcal{H}$.
	\[\xymatrix@C=2.3cm{
			\Delta_{\EqHyp}(X) \ar@{.>}[d]_{\Delta_\EqHyp(f)} \ar[dr]^{h} & \\
			\Delta_\EqHyp(U_{\EqHyp}(\mathcal{H})) \ar[r]_{\epsilon_{\mathcal{H}}} & \mathcal{H}
	}\]
	Where $\Delta_\EqHyp(U_{\EqHyp}(\mathcal{H})) = (\emptyset, V_\mathcal{H}, \{\bullet\}, ?_{V_\mathcal{H}}, ?_{V_\mathcal{H}}, !_{V_\mathcal{H}})$ and $\epsilon_{\mathcal{H}} = (?_{E_\mathcal{H}}, \id{V_\mathcal{H}}, g)$.
	Note that, since $\Delta_{\EqHyp}(X)$ has the empty set as object of edges, $h_E = ?_{E_\mathcal{H}}$, then, the unique arrow that fits in the diagram is $\Delta_{\EqHyp}(f) = (?_{E_\mathcal{H}}, h_V, \id{\{\bullet\}})$.
\end{proof}

We now define another functor $T: \EqHyp \to \hyp$, which ``forgets'' the quotient part, mapping each hypergraph with equivalence $\mathcal{G} = (E_\mathcal{G}, V_{\mathcal{G}}, C_\mathcal{G}, s_\mathcal{G}, t_\mathcal{G}, q_\mathcal{G})$ onto $T(\mathcal{G})=(E_{\mathcal{G}}, V_{\mathcal{G}}, s_\mathcal{G}, t_{\mathcal{G}})$. Then, we have the following result.

\begin{proposition}
	$T$ has a left adjoint $L: \hyp \to \EqHyp$.
\end{proposition}

\begin{proof}
	Let $\mathcal{G}$ be a hypergraph, and define $L(\mathcal{G}) := (E_\mathcal{G}, V_{\mathcal{G}}, \{\bullet\}, s_\mathcal{G}, t_\mathcal{G}, !_{V_\mathcal{G}})$. Let now $h: L(\mathcal{G})\to \mathcal{H}$ be a morphism in $\EqHyp$, and consider the following situation
	\[\xymatrix@C=2.3cm{
		L(\mathcal{G}) \ar@{.>}[d]_{L(f)} \ar[dr]^{h}\\ L(T(\mathcal{H})) \ar[r]_{\epsilon_{\mathcal{H}}} & \mathcal{H}
	}\]
	for $L(T(\mathcal{H}))=(E_\mathcal{H}, V_{\mathcal{H}}, \{\bullet\}, s_\mathcal{H}, t_\mathcal{H}, !_{V_\mathcal{H}})$. Then, $\epsilon_\mathcal{H} = (\id{E_\mathcal{H}}, \id{V_\mathcal{H}}, h_C)$ (by \Cref{rem:eqhyp_morphs}, the last component is uniquely determined by the first two), and $L(f)$ must be $(h_E, h_V, \id{\{\bullet\}})$.
\end{proof}

\begin{remark}
	$T$ is faithful. Indeed, consider two morphisms $h = (h_E, h_V, h_C)$ and $k = (k_E, k_V, k_C)$, and suppose $T(h) = T(k)$, that is, $(h_E, h_V) = (k_E, k_V)$.
	By \Cref{rem:eqhyp_morphs}, we can conclude also $h_C = k_C$, and hence the faithfulness of $T$.
\end{remark}

Let now $K: \EqHyp \to \Set$ be the functor which sends each hypergraph with equivalence $\mathcal{G} = (E, V, C, s, t, q)$ onto $K(\mathcal{G}) = C$, and each morphism $(h_E, h_V, h_C)$ to $h_C$.

\begin{proposition}\label{prop:eqhyp_complete}
	\EqHyp \, is complete and cocomplete, and $T$ preserves limits and colimits.
\end{proposition}

\begin{proof}
	Let $D: \cat{I}\to \EqHyp$ be a diagram, and, for each $i \in \cat{I}$, $D(i) = (E_i, V_i, C_i, s_i, t_i, q_i)$.
	Suppose now $(E, V, s, t)$, together with morphisms $(\pi_i^E, \pi_i^V)$, be the limit of $T \circ D$.
	Then, $V$, together with $(q_i\circ \pi_V^i)_{i\in \cat I}$, is a cone for $K \circ D$. Indeed, let $\alpha: i \to j$ be an arrow of $\cat I$, $D(\alpha) = (h_E,h_V, h_C)$.
	By definition of $T$, $(T \circ D)(\alpha) = (h_E, h_V)$, hence we have
	\[\xymatrix{
		& V \ar[dl]_{\pi_V^i} \ar[dr]^{\pi_V^j} & \\V_i\ar[rr]^{h_V}\ar[d]_{q_i}&&V_j\ar[d]^{q_j}\\C_i\ar[rr]_{h_C}&&C_j
	}\]
	Suppose now that $L$, with morphisms $(l_i)_{i\in \cat I}$ be the limit of $K\circ D$. Hence, we have an arrow $l:V \to L$, which is not epi in general.
	Let then $l = m \circ q$ be the epi-mono factorization of it. Consider the following situation, where the outer rectangle commutes by definition, and the dotted arrow is yielded by ({\color{red}cite left lifting prop}).
	\[\xymatrix{
			V\ar[r]^{\pi_V^i}\ar[d]_{q}&V_i\ar[r]^{q_i}&C_i\ar[d]^{\id{C_i}}\\C\ar[r]_{m}\ar@{.>}[urr]^{\pi_C^i}&L\ar[r]_{l_i}&{C_i}
	}\]
	Thus, $(E, V, C, s, t, q)$, together with $(\pi_E^i, \pi_V^i, \pi_C^i)$ is a cone over $D$. {\color{red} remain to show that this cone is terminal}
	
	Suppose now $(E', V', s', t')$, together with $(\kappa_E^i, \kappa_V^i)_{i \in \cat I}$, be the colimit of $T \circ D$, and $C'$, with $(c_i)_{i \in \cat I}$ be the colimit of $K \circ D$.
	Then, we have the folliwing situation.
	\[\xymatrix{
								  &  {V'} &                             \\
		V_i \ar[ur]^{\kappa_i} \ar[rr]^{h_V} \ar[d]_{q_i} && V_j \ar[ul]_{\kappa_j} \ar[d]^{q_j}\\
		C_i \ar[dr]_{c_i}      \ar[rr]_{h_C}              &&C_j  \ar[dl]^{c_j}                  \\
								  &  {C'} &
	}\]
	Then, $C'$ with morphisms $(c_i \circ q_i)_{i\in \cat I}$ is a conone for $U \circ D$. Then, there exists a unique morphism $q': V' \to C'$ such that $q'\circ \kappa_V^i = c_i\circ q_i$.
	Such morphism is epi ({\color{red}cite Lemma 1.3.45 of the thesis}), and thus $(E', V', C', s', t', q')$, together with $(\kappa_E^i, \kappa_V^i, c_i)_{i\in \cat I}$ is the colimit of $D$.
	
\end{proof}

\begin{corollary}\label{cor:mono1}
	Let $h = (h_E, h_V, h_C): \mathcal{G \to H}$ be an arrow in  $\EqHyp$. Then it is a mono if and only if $T(h)$ is a mono.
\end{corollary}

\begin{proof}
	The ``if'' part is given by the faithfulness of $T$.
	The ``only if'' part is given by \Cref{rem:eqhyp_morphs}.
\end{proof}

\begin{corollary}\label{cor:mono2}
	If $h = (h_E, h_V, h_C): \mathcal{G\to H}$ is a regular mono in $\EqHyp$, then $h_E, h_V$ and $h_C$ are all monos.
\end{corollary}

\begin{proof}
	If $h$ is mono, from \Cref{cor:mono1} we have that $h_E$ and $h_V$ are monos. Suppose now $f, g: \mathcal{H\rightrightarrows K}$ be the arrows equalized by $h$. Then, we have:
	\begin{align*} f_C \circ h_C \circ q_\mathcal{G} &= f_C\circ q_\mathcal{H} \circ h_V\\&=q_\mathcal{K}\circ g_V \circ h_V\\&=q_\mathcal{K}\circ f_V \circ h_V\\&=g_C \circ h_C \circ q_\mathcal{G} \end{align*}
	Since $q_\mathcal{G}$ is epi, we have $f_C \circ h_C = g_C \circ h_C$, hence $h_C$ is an equalizer for $f_C$ and $g_C$, and thus a mono.
\end{proof}

\begin{proposition}
	Let $h = (h_E, h_V, h_C): \mathcal{G \to H}$ be a regular mono in $\EqHyp$.
	Then, $h_E$ and $h_V$ are monos and $(K, \pi_1, \pi_2)$ is the kernel pair of $q_\mathcal{H}\circ h_V$ if and only if $(K, \pi_1, \pi_2)$ is the kernel pair of $q_\mathcal{G}$.
\end{proposition}

\begin{proof}
	By \Cref{cor:mono2}, we have that $h_E, h_V$ and $h_C$ are all monos.
	Hence, by \Cref{cor:kermono}, $(K, \pi_1, \pi_2)$ is the kernel pair of $q_\mathcal{G}$ if and only if it is the kernel pair also of $h_C \circ q_\mathcal{G}$, since $h_C$ is mono by hypothesis.
	The thesis follows from $h_C \circ q_\mathcal{G} = q_H \circ h_V$, and from the hypothesis of $h_E$ mono.
\end{proof}


\begin{remark}
    It is possible to restate the last proposition, by \Cref{ex:kernel_pairs_in_Set}, as 
    \begin{displayquote}
    \textit{$h_E$ and $h_V$ are mono and, for every $v, v'\in V_H$, $q_H(h_V(v))=q_H(h_V(v'))$ if and only if $q_G(v)=q_G(v')$}
    \end{displayquote}
    That is, a regular mono in $\EqHyp$ is a morphism that both reflects and preserves equivalences.
    % besides preserving them.
\end{remark}

Let us turn to another functor $\EqHyp \to \hyp$.

\begin{definition}
The \emph{quotient functor} $Q:\EqHyp\to \hyp $ is defined as
the one sending $(E, V, C, s, t, q)$ to $(E, C, q^{\star}\circ s, q^{\star}\circ t)$ and an arrow $(h_E, h_V, h_C)$ to $(h_E, h_C)$.
\end{definition}

\begin{remark}
	The action of the functor on a morphism of hypergraphs with equivalences gives a morphism of hypergraphs,
	in fact $q^{\star}_\mathcal{H} \circ s_\mathcal{H} \circ h_E = q^{\star}_\mathcal{H} \circ h_V^\star \circ s_\mathcal{G} = h_C^\star \circ q^{\star}_\mathcal{G} \circ s_\mathcal{G}$.
	The same is valid for $t_\mathcal{H}$ and $t_\mathcal{G}$. 
\end{remark}

\begin{lemma}\label{lemma:quot_funct_left_adj}
    $Q$ is a left adjoint.
\end{lemma}

\begin{proof}
	Let $R((A, B, s, t))$ be $(A, B, B, s, t, \id{B})$, so that $Q(R((A, B, s, t))) = (A, B, s, t)$. Now, suppose that $h = (h_E, h_V): Q((E, V, C, s', t', q)) \to (A, B, s, t)$  is an arrow in $\hyp$, and consider the triple $(h_E, h_V, h_V \circ q)$. Since $h$ is a morphism of $\hyp$, we have $h_V^{\star}\circ q^{\star}\circ s'= s\circ h_E$ and $h_V^{\star}\circ q^{\star}\circ t' = t\circ h_E$.
    Then we have the following squares
    \[\xymatrix{
		    E \ar[r]^{h_E}\ar[d]_{s'} & A \ar[d]^{s} & E \ar[r]^{h_E}\ar[d]_{t'} & A \ar[d]^{t} & V \ar[r]^{h_V \circ q} \ar[d]_{q} & B \ar[d]^{\id{B}} \\
		    V^{\star} \ar[r]_{h_V^\star\circ q^\star} & B^{\star} & V^{\star} \ar[r]_{h_V^\star\circ q^\star} & B^{\star} & C \ar[r]_{h_V} & B
    }\]

    \iffalse\[
        \begin{tikzcd}
            E \arrow[r, "{h_E}"] \arrow[d, "{s_G}"swap] & A \arrow[d, "s"] \\
            V \arrow[r, "{h_V \circ q}"swap] & B
        \end{tikzcd}
        \qquad
        \begin{tikzcd}
            E \arrow[r, "{h_E}"] \arrow[d, "{t_G}"swap] & A \arrow[d, "t"] \\
            V \arrow[r, "{h_V \circ q}"swap] & B
        \end{tikzcd}
        \qquad
        \begin{tikzcd}
            V \arrow[r, "{h_V\circ q}"] \arrow[d, "q" swap] & B \arrow[d, "{id_B}"] \\
            C \arrow[r, "{h_V}"swap] & B
        \end{tikzcd}
    \]\fi
    We have therefore found a morphism $(E, V, C, s', t', q) \to R((A, B, s, t))$ whose image through $Q$ fits in the diagram below
    \[\xymatrix{
		    (A, B, s, t) \ar[r]^{(\id{A}, \id{B})} & (A, B, s, t)\\
		    (E, C, q^\star\circ s', q^\star\circ t') \ar[u]^{Q(h_E, h_V \circ q, h_V)} \ar[ur]_{(h_E, h_V)}
    }\]   
    Such arrow is unique. Suppose $f = (f_E, f_V, f_C)$ to be another arrow with such property.
    Then, it must be $(\id A, \id B) \circ Q(f) = (f_E, f_C) = (h_E, h_C)$. Finally, $f_C = f_V \circ q = h_V \circ q$. 
\end{proof}

\begin{proposition}\label{prop:quot_creat_colims}
    $Q$ creates colimits.
\end{proposition}


\begin{proof}
    Since $Q$ is a left adjoint, it preserves colimits.
	Let $D: \cat I \to \EqHyp$ be a diagram, and let $\mathcal{C}$, together with  $(c_i)_{i\in \cat I}$ be the colimit of $Q \circ D$, where $\mathcal{C} = (A, C, q\circ s, q\circ t)$, and $D(i)$ is $(A_i, B_i, C_i, s_i, t_i, q_i)$.
	\iffalse
	Let now $T: \EqHyp \to \Set$ be the functor mapping each graph with equivalence onto its second component, $T((X, Y, Z, x, y, z)) = Y$, and each morphims onto its second component.\fi
	Let $((\kappa_i)_{i\in \cat I}, V)$ be the colimit of $U_\EqHyp \circ D$.
	Consider the following situation
	\[\xymatrix{
								  &  {V} &                             \\
		V_i \ar[ur]^{\kappa_i} \ar[rr]^{h_V} \ar[d]_{q_i} && V_j \ar[ul]_{\kappa_j} \ar[d]^{q_j}\\
		C_i \ar[dr]_{c_i}      \ar[rr]_{h_C}              &&C_j  \ar[dl]^{c_j}                  \\
								  &  {C} &
	}\]
	Now, since $((c^i_C \circ q_i)_{i \in \cat I}, C)$ is a cocone for $U_\EqHyp \circ D$, there exists a unique $q: V \to C$, which is epi by \Cref{lemma:nat_trans_reg_epi_canonical_arrow_reg_epi}.
	Consider now the functor $W: \EqHyp \to \Set$ mapping each $(X, Y, Z, x, y, z)$ onto $X$, and each morphism on its first component.
	By \Cref{prop:eqhyp_complete} and \Cref{lemma:limits_of_presheaves}, we have that $((c_E^i)_{i \in \cat I}, E)$ is the colimit of $W \circ D$.
	Notice that $((\kappa_i \circ s_i)_{i \in \cat{I}}, B)$ and $((\kappa_i \circ t_i)_{i \in \cat I}, B)$ are cocones for $W \circ D$, so let $s$ and $t$ be, respectively, the mediating arrow for the first one and the mediating arrow for the second one. It remains now to show that $(E, V, C, s, t, q)$, together with $(c_E^i, \kappa_i, c_C^i)_{i \in \cat I}$, is a colimit for $D$, but this follows by the proof of \Cref{prop:eqgrph_complete}.
\end{proof}

\begin{example}
	$Q$ does not preserve limits. Indeed, let $\mathcal{G}_1 = (E_1, A, A, s_1, t_1, id_A)$, $\mathcal{G}_2 = (E_2, B, B, s_2, t_2, id_B)$ and $\mathcal{G}_3 = (E_3, A + B, \terminal, s_3, t_3, !_{A + B})$, and let $h = (h_E, \iota_A, !_A): \mathcal{G}_1 \to \mathcal{G}_3$, $k = (k_E, \iota_B, !_B): \mathcal{G}_2 \to \mathcal{G}_3$, where $(\iota_A, \iota_B, A + B)$ is the coproduct of $A$ and $B$, $\terminal$ is the intial object (in $\Set$, the singleton set as shown in \Cref{ex:set_init_term}), and $!_X$ the unique arrow $X \to \terminal$.
	The following two diagrams show the pullback of $h$ and $k$ and the pullback of $Q(h)$ and $Q(k)$, on the second component (the vertices of the graphs)
	\[\xymatrix{
			\initial \ar[r]^{p_1} \ar[d]_{p_2}& A \ar[d]^{\iota_A} \\
			B \ar[r]_{\iota_B} & A+B
	}
	\qquad
	\xymatrix{
		{A \times B} \ar[r]^{\pi_A} \ar[d]_{\pi_B} & A \ar[d]^{!_A} \\
		B \ar[r]_{!_B} & \terminal
	}\]

	But the arrow $\initial \to A \times B$ is not an epi in general (this is easy to see taking $\Set$ as example), hence such pullback is not preserved by $Q$.
\end{example}

\begin{lemma}\label{lemma:stab}
	In $\EqHyp$, pushouts along regular monos are stable.
\end{lemma}

\begin{proof}
	Let $\mathcal{G}_i = (A_i, B_i, C_i, s_i, t_i, q_i)$, $\mathcal{G}'_i=(A'_i, B'_i, C'_i, s'_i, t'_i, q'_i)$, for $i \in \{1, 2, 3, 4\}$, be hypergraphs with equivalence, 
	and, in the diagram below, suppose all the vertical faces are pullbacks, the bottom face is a pushout and $h$ is regular mono.
                \[
		\xymatrix@C=10pt@R=6pt{&\mathcal{G}_1'\ar[dd]|\hole_(.65){a}\ar[rr]^{h'} \ar[dl]_{k'} && \mathcal{G}_2' \ar[dd]^{b} \ar[dl]_{t'} \\ \mathcal{G}_3'  \ar[dd]_{c}\ar[rr]^(.7){p'} & & \mathcal{G}_4' \ar[dd]_(.3){d}\\&\mathcal{G}_1\ar[rr]|\hole^(.65){h} \ar[dl]^{k} && \mathcal{G}_2 \ar[dl]^{t} \\\mathcal{G}_3 \ar[rr]_{p} & & \mathcal{G}_4 }
	\]
	By \Cref{prop:eqhyp_complete} and \Cref{cor:mono1}, the following cubes in $\Set$ have pushouts as bottom faces and pullbacks as vertical faces, hence their top faces are pushouts.
	\[
        \xymatrix@C=10pt@R=6pt{&A_1'\ar[dd]|\hole_(.65){a_E}\ar[rr]^{h_E'} \ar[dl]_{k_E'} && A_2' \ar[dd]^{b_E} \ar[dl]_{t_E'} \\ A_3'  \ar[dd]_{c_E}\ar[rr]^(.7){p_E'} & & A_4' \ar[dd]_(.3){d_E}\\&A_1\ar[rr]|\hole^(.65){h_E} \ar[dl]^{k_E} && A_2 \ar[dl]^{t_E} \\A_3 \ar[rr]_{p_E} & & A_4 }
	\qquad
        \xymatrix@C=10pt@R=6pt{&B_1'\ar[dd]|\hole_(.65){a_V}\ar[rr]^{h_V'} \ar[dl]_{k_V'} && B_2' \ar[dd]^{b_V} \ar[dl]_{t_V'} \\ B_3'  \ar[dd]_{c_V}\ar[rr]^(.7){p_V'} & & B_4' \ar[dd]_(.3){d_V}\\&B_1\ar[rr]|\hole^(.65){h_V} \ar[dl]^{k_V} && B_2 \ar[dl]^{t_V} \\B_3 \ar[rr]_{p_V} & & B_4 }
	\]
	Consider now the following pullbacks.
	\[\xymatrix{
			Y \ar[d]_{y_2}\ar[r]^{y_1}&{C_4'}\ar[d]^{d_C}&U\ar[d]_{u_2}\ar[r]^{u_1}&{C_4'}\ar[d]^{d_C}&T\ar[r]^{x_1}\ar[d]_{x_2}&U\ar[d]^{u_2}\\
			C_3\ar[r]_{p_C}&C_4&C_2\ar[r]_{t_C}&C_4&C_1\ar[r]_{h_C}&C_2
	}\]
	Thus, [???]\todo{inserire in sezione 1 la proposizione} yields the following situation, in which the bottom face is a pushout, and the vertical faces are pullbacks, hence the top face is a pushout too.
        \[\xymatrix@C=10pt@R=6pt{&T\ar[dd]|\hole_(.65){x_2}\ar[rr]^{x_1} \ar@{.>}[dl]_{w} && U \ar[dd]^{u_2} \ar[dl]_{u_1} \\ Y  \ar[dd]_{y_2}\ar[rr]^(.7){y_1} & & C_4' \ar[dd]_(.3){d_C}\\&C_1\ar[rr]|\hole^(.65){h_C} \ar[dl]^{k_C} && C_2 \ar[dl]^{t_C} \\C_3 \ar[rr]_{p_C} & & C_4 }\]
	By the proof of \Cref{prop:eqhyp_complete}, we have that $m_2 \circ q'_3 : B'_3 \to Y$ and $m_3 \circ q'_2 : B_2' \to U$ are two epi-mono factorizations, with $m_2$ and $m_3$ monos.
	At the same way, let the following square to be a pullback.
	\[\xymatrix{S\ar[r]^{s_1}\ar[d]_{s_2}&C'_2\ar[d]^{m_3}\\T\ar[r]_{x_1}&U}\]
	Hence, in the following diagram, the outer rectangle is a pullback.
	\[\xymatrix{
			S\ar[r]^{s_1}\ar[d]_{s_2}&C'_2\ar[d]^{m_3}\\T\ar[r]_{x_1}\ar[d]_{x_2}&U\ar[d]^{u_2}\\C_1\ar[r]_{h_C}&C_2
	}\]
	By the same argument as before, there exists a mono $m_1$ such that $m_1\circ q_1' : B_1' \to S$.
	% Uguaglianze da capire se servono davvero

	We have to show that the top face of the cube ate the beginning of the proof is a pushout.
	Suppose then that $z: \mathcal{G}_2' \to \mathcal{H}$ and $w: \mathcal{G}_3' \to \mathcal{H}$, with $\mathcal{H} = (E, V, C, s, t, q)$,
	are two morphisms such that $z \circ h' = w \circ h'$, and let $v_V: B'_4 \to V$ the arrow induced by $z_V$ and $w_V$.
	We want to construct the dotted arrow $v_C$ which fits in the diagram below.
	\[
	\xymatrix@C=10pt@R=6pt{
		&B_1'\ar[dd]|\hole_(.65){q_1'}\ar[rr]^{h_V'} \ar[dl]_{k_V'} && B_2' \ar[dd]|\hole_(.65){q_2'} \ar[dl]_{t_V'}\ar[dr]^{z_V}\\
		B_3'  \ar[dd]_{q_3'}\ar[rr]^(.7){p_V'} & & B_4' \ar[dd]_(.3){q_4'} \ar[rr]^(.7){v_V}&&  V\ar[dd]^{q}\\
						       &C_1'\ar[rr]|\hole^(.65){h_C'} \ar[dl]_{k_C'} && C_2' \ar[dl]^{t_C'}\ar[dr]^{z_C}\\
		C'_3 \ar[rr]_{p_C'} & & C'_4\ar@{.>}[rr]_{v_C} && C
	}\]
	By \Cref{lem:mpo}, we know that the top face of the cube below is a pushout.
\[\xymatrix@C=15pt@R=15pt{&K_{s_2\circ m_1\circ q'_1}\ar[dd]|\hole_(.65){\pi_{s_2\circ m_1\circ q'_1}^{1}}\ar[rr]^{k_{h'_2}} \ar[dl]_{k_{k'_2}} && K_{m_2\circ q'_2} \ar[dd]^{\pi_{m_2\circ q'_2}^{1}} \ar[dl]_{k_{t'_2}} \\ K_{m_3\circ q'_3}  \ar[dd]_{\pi_{m_3\circ q'_3}^{1}}\ar[rr]^(.7){k_{p'_2}} & & K_{q'_4} \ar[dd]_(.3){\pi_{q'_4}^{1}}\\&B'_1\ar[rr]|\hole^(.65){h'_2} \ar[dl]^{k'_2} && B'_2 \ar[dl]^{t'_2} \\B'_3 \ar[rr]_{p'_2} & & C_4}\]
	And, since $m_3$ and $m_2$ are monos,
	\[\begin{split}
		q'_3 \circ \pi_{m_3 \circ q'_3}^1 = q'_3 \circ \pi_{m_3\circ q_3'}^2
	\end{split}\qquad\begin{split}
		q_2' \circ \pi_{m_2 \circ q_2'}^1 = q'_2 \circ \pi_{m_2 \circ q_2'}^2
	\end{split}\]
	Computing, we obtain
	\[\begin{split}
		q \circ v_V \circ \pi_{q_4'}^1 \circ k_{p_V'} &= q \circ v_V \circ p_V' \circ \pi_{m_3 \circ q_3'}^1 \\
							      &= q \circ w_V \circ \pi_{m_3 \circ q_3'}^1 \\
							      &= w_C \circ q_3' \circ \pi_{m_3 \circ q_3'}^1 \\
							      &= w_C \circ q_3' \circ \pi_{m_3 \circ q_3}^2 \\
							      &= q \circ w_V \circ \pi_{m_3 \circ q_3'}^2 \\
							      &= q \circ v_V \circ p_V' \circ \pi_{m_3 \circ q_3'}^2 \\
							      &= q \circ v_V \circ \pi_{q'_4}^2 \circ k_{p'_V}
	\end{split}\qquad\begin{split}
		q \circ v_V \circ \pi_{q'_4}^1 \circ k_{t_V'} &= q \circ v_V \circ t_V' \circ \pi_{m_2 \circ q_2'}^1 \\
							      &= q \circ z_V \circ \pi_{m_2 \circ q_2'}^1 \\
							      &= z_C \circ q_2' \circ \pi_{m_2 \circ q_2'}^1 \\
							      &= z_C \circ q_2' \circ \pi_{m_2 \circ q_2}^2 \\
							      &= q \circ z_V \circ \pi_{m_2 \circ q_2'}^2 \\
							      &= q \circ v_V \circ t_V' \circ \pi_{m_2 \circ q_2'}^2 \\
							      &= q \circ v_V \circ \pi_{q'_4}^2 \circ k_{t'_V}
	\end{split}\]
	Since the previous cube has a pushout as top face, by universal property, we have
	\[
		q \circ v_V \circ \pi_{q'_4}^1 = q \circ v_V \circ \pi_{q'_4}^2
	\]
	hence, $v_C$ is the mediating arrow.
	\[
		v_C \circ q'_4 \circ \pi_{q'_4}^1 = v_C \circ q'_4 \circ \pi_{q'_4}^2
	\]
\end{proof}

\begin{lemma}\label{lemma:van_kampen}
	In $\EqHyp$, pushouts along regular monos are $\reg(\EqHyp)$-Van Kampen.
\end{lemma}

\begin{proof}
	In lieu of \Cref{lemma:stab}, it is enough to proof that, given a cube as the one below, with pullbacks as back faces, pushouts as bottom and top faces and such that $h$ is a regular mono,
	the front faces are pullbacks too, where $\mathcal{G}_i = (A_i, B_i, C_i, s_i, t_i, q_i)$, $\mathcal{G}'=(A_i', B_i', C_i', s_i', t_i', q_i')$, for $i = 1, 2, 3, 4$.
  	\[
	\xymatrix@C=10pt@R=6pt{&\mathcal{G}_1'\ar[dd]|\hole_(.65){a}\ar[rr]^{h'} \ar[dl]_{k'} && \mathcal{G}_2' \ar[dd]^{b} \ar[dl]_{t'} \\ \mathcal{G}_3'  \ar[dd]_{c}\ar[rr]^(.7){p'} & & \mathcal{G}_4' \ar[dd]_(.3){d}\\&\mathcal{G}_1\ar[rr]|\hole^(.65){h} \ar[dl]^{k} && \mathcal{G}_2 \ar[dl]^{t} \\\mathcal{G}_3 \ar[rr]_{p} & & \mathcal{G}_4 }
	\]
	By \Cref{prop:eqhyp_complete} and \Cref{cor:mono_in_EqGrph}, the following two cubes have $\mathcal{M}$-pushouts as bottom faces and pullbacks as back faces,
	thus their front faces are pullbacks too.
	\[
        \xymatrix@C=10pt@R=6pt{&A_1'\ar[dd]|\hole_(.65){a_E}\ar[rr]^{h_E'} \ar[dl]_{k_E'} && A_2' \ar[dd]^{b_E} \ar[dl]_{t_E'} \\ A_3'  \ar[dd]_{c_E}\ar[rr]^(.7){p_E'} & & A_4' \ar[dd]_(.3){d_E}\\&A_1\ar[rr]|\hole^(.65){h_E} \ar[dl]^{k_E} && A_2 \ar[dl]^{t_E} \\A_3 \ar[rr]_{p_E} & & A_4 }
	\qquad
        \xymatrix@C=10pt@R=6pt{&B_1'\ar[dd]|\hole_(.65){a_V}\ar[rr]^{h_V'} \ar[dl]_{k_V'} && B_2' \ar[dd]^{b_V} \ar[dl]_{t_V'} \\ B_3'  \ar[dd]_{c_V}\ar[rr]^(.7){p_V'} & & B_4' \ar[dd]_(.3){d_V}\\&B_1\ar[rr]|\hole^(.65){h_V} \ar[dl]^{k_V} && B_2 \ar[dl]^{t_V} \\B_3 \ar[rr]_{p_V} & & B_4 }
	\]
	On the other hand we can consider the diagrams below, in which the inner squares are pullbacks.
	Since the outer diagrams commute, by definition of morphism of $\EqHyp$, then we have the existence of $m_2\colon C'_2\to U$, $m_3\colon C'_3\to Y $, $a_3\colon B'_3\to Y$ and $a_2\colon B'_2\to Y$.
	\[\xymatrix{C'_3 \ar@{.>}[dr]^{m_3} \ar@/^.3cm/[drr]^{p'_3} \ar@/_.3cm/[ddr]_{d_3} &&& C'_2 \ar@{.>}[dr]^{m_2} \ar@/^.3cm/[drr]^{t'_3} \ar@/_.3cm/[ddr]_{d_2}\\&Y \ar[r]_{y_1} \ar[d]_{y_2}& C'_4\ar[d]^{d_3}&& U \ar[d]_{u_2}\ar[r]_{u_1}& C'_4 \ar[d]^{d_3}\\&C_3 \ar[r]_{p_3}& C_4 &&C_2 \ar[r]_{t_3}& C_4 }\]

		\[\xymatrix{B'_3 \ar@{.>}[dr]^{a_3} \ar[r]^{p'_2} \ar[d]_{q'_3} &B'_4\ar[dr]^{q'_4}&& B'_2\ar@{.>}[dr]^{a_2} \ar[r]^{t'_2} \ar[d]_{q'_2} &B'_4\ar[dr]^{q'_4}\\C'_3 \ar[dr]_{d_3} &Y \ar[r]_{y_1} \ar[d]_{y_2}& C'_4\ar[d]^{d_3}&C_2' \ar[dr]_{d_2}& U \ar[d]_{u_2}\ar[r]_{u_1}& C'_4 \ar[d]^{d_3}\\&C_3 \ar[r]_{p_3}& C_4 &&C_2 \ar[r]_{t_3}& C_4 }\]

	Now, notice that  $m_3$ and $m_2$ are monos because $d_3$ and $d_2$ are regular monos. By the proof of \Cref{prop:eqhyp_complete}, to conclude it is enough to show that
	\[m_3\circ q'_3 = a_3 \qquad m_2\circ q'_2=a_2\]

	Indeed, if the previous equations hold, then $C'_3$ and $C'_2$ are epi-mono factorizations of $a_3$ and $a_2$ and the thesis follows from \Cref{cor:unique} and the proof of \Cref{prop:eqhyp_complete}.

	No if we compute we have:
	\[\begin{split}
		y_1\circ a_3&= q'_4\circ p'_2\\&=p'_3 \circ q'_3\\&=y_1\circ m_3\circ q'_3  
	\end{split}\qquad \begin{split}
		u_1\circ a_2&= q'_4\circ t'_2\\&=t'_3 \circ q'_3\\&=u_1\circ m_2\circ q'_2  
	\end{split}\]
	\[\begin{split}
	y_2\circ a_3&= d_3\circ q_3'\\&=y_2\circ m_3\circ q'_3
	\end{split}\qquad \begin{split}
	u_2\circ a_2&= d_2\circ q_2'\\&=u_2\circ m_2\circ q'_2
	\end{split}\]
	And we have done.
\end{proof} 



\subsection{title}

\section{EGGS}

\section{Conclusions and further works}
\todo{Some very nice conclusions}
\bibliography{biblio}

\appendix
\section{Some properties of comma categories}
In this section we will briefly recall the definition of the comma category \cite{mac2013categories} associated to two functors and some of its properties.
\begin{definition}\index{category!comma -}
	Let $L\colon \A\to \X$ and  $R\colon \B\rightarrow \X$ be two functors with the same codomain, the \emph{comma category} $\comma{L}{R}$ is the category in which
	\begin{itemize}
		\item objects are triples $(A, B, f)$ with $A\in \A$, $B\in \B$, and $f\colon L(A)\rightarrow R(B)$; 
		\item a morphism $(A, B, f)\rightarrow (A', B', g)$ is a pair $(h, k)$ with $h\colon A\rightarrow A'$ in $\A$ and $k\colon B\rightarrow B'$ in $\B$ such that the following diagram commutes
		\[\xymatrix@C=25pt{L(A) \ar[r]^{L(h)} \ar[d]_{f} & L(A') \ar[d]^{g}\\ R(B) \ar[r]_{R(k)}& R(B')}\]
	\end{itemize}
\end{definition} 
We have two forgetful functors 	$U_L\colon \comma{L}{R}\to \A$ and $U_R\colon \comma{L}{R} \to \B$ given, respectively by
\[
\begin{split}
	\functor[l]{(A,B, f)}{(h,k)}{(A', B', g)}
	& \functormapsto
	\rfunctor{A}{h}{A'}
\end{split}\quad 
\begin{split}
	\functor[l]{(A,B, f)}{(h,k)}{(A', B', g)}
	& \functormapsto
	\rfunctor{B}{k}{B'}
\end{split}
\]


Given $L\colon \A\to \X$ and $R\colon \B\to \X$, we can also consider their duals $L^{op}\colon \A^{op}\to \X^{op}$ and $R^{op}\colon \B^{op}\to \X^{op}$.  An arrow $f\colon L(A)\to R(B)$ in $\X$ is the same ting as an arrow $f\colon R^{op}(B)\to L^{op}(A)$ in $\X^{op}$, thus $\left(\comma{L}{R}\right)$ and $\comma{R^{op}}{L^{op}}$ have the same objects. Moreover, the commutativity in  $\X$ of the square
\[\xymatrix@C=25pt{L(A) \ar[r]^{L(h)} \ar[d]_{f} & L(A') \ar[d]^{g}\\ R(B) \ar[r]_{R(k)}& R(B')}\]
is tantamount to the commutativity in $\X^{op}$ of the square 
\[\xymatrix@C=25pt{ R(B') \ar[d]_{g} \ar[r]^{R(k)}& R(B) \ar[d]^{f}\\ L(A') \ar[r]_{L(h)}  & L(A) }\]
Summing up we have just proved the following fact.
\begin{proposition}\label[proposition]{prop:dual}
	$(\comma{L}{R})^{op}$ is equal to	$\comma{R^{op}}{L^{op}}$, moreover $U^{op}_L=U_{L^{op}}$ and $U^{op}_R=U_{R^{op}}$.
\end{proposition}


\begin{lemma}\label[lemma]{colim}
	Let $L\colon \A\to \X$ and $R\colon \B\to \X$ be functors and $F\colon \D\to\comma{L}{R}$ be a diagram such that $L$ preserves colimits along $U_L\circ F$. Then the family $\{U_L, U_R\}$ jointly creates colimits of $F$ (see \cite{castelnovo2023thesis,castelnovo2022newcriterionmathcalmmathcalnadhesivity}).
\end{lemma}
\begin{proof}
	Suppose that $U_L\circ F$ and $U_R\circ F$ have colimiting cocones  $\left(A, \{a_D\}_{D\in \D}\right)$ and $\left(B, \{b_D\}_{D\in \D}\right)$ respectively. By hypothesis $\left(L(A), \{L\left(a_D\right)\}_{D\in \D}\right)$ is colimiting for $L\circ U_L\circ F$. Now, if we define
	\[F(D):=(A_D, B_D, f_D)\]
	then we have arrows $R(a_i)\circ f_D\colon L(A_D)\to R(B)$ that forms a cocone on $L\circ U_L\circ F$: if $d\colon D\to D'$ is an arrow in $\D$ then $F(d)$ is an arrow in $\comma{L}{R}$ and so
	\begin{align*}
		R\left(b_{D'}\right)\circ f_{D'}\circ L(U_L(F(d)))&=R\left(b_{D'}\right)\circ R\left(U_R\left(F(d)\right)\right)\circ f_D\\&=R\left(b_{D'}\circ U_R\left(F(d)\right)\right)\circ f_D\\&=R\left(b_D\right)\circ f_D
	\end{align*}
	Thus there exists $f\colon L(A)\rightarrow R(B)$ such that
	\[\xymatrix@C=35pt{ L(A_D) \ar[r]^{L(a_D)} \ar[d]_{f_D}& L(A) \ar[d]^{f}\\ R(B_D) \ar[r]_{R(b_D)} &R(B)}\]
	Notice that $f$ is the unique arrow in $\X$ wich makes $\left(a_D, b_D\right)$ an arrow $\left(A_D, B_D, f_D\right)\to \left(A, B, f\right)$ of $\comma{L}{R}$. If we show that $\left((A, B, f), \left\{(a_D, b_D)\right\}_{D\in \D}\right)$ is colimiting for $F$ we are done.
	
	First of all, let us show that it is a cocone. Given $d\colon D\to D'$ in $\D$ we have:
	\begin{align*}
		\left(a_{D'}, b_{D'}\right)\circ F(d)&=	\left(a_{D'}, b_{D'}\right)\circ \left(U_L(F(d)), U_R(F(d))\right)\\&=\left( a_{D'}\circ U_L(F(d)),  b_{D'}\circ U_R(F(d)) \right)\\&=\left(a_D, b_D\right)
	\end{align*}
	For the colimiting property, let $\left((X, Y, g), \left\{\left(x_D, y_D\right)\right\}_{D\in \D}\right)$ be another cocone on $F$. In particular $\left(X, \left\{x_D\right\}_{D\in \D}\right)$ and $\left(Y, \left\{y_D\right\}_{D\in \D}\right)$ are cocones on $U_L\circ F$ and $U_R\circ F$ respectively, so we have uniquely determined arrows $x\colon A\rightarrow X$ and $y\colon B\rightarrow Y$ such that 
	\[x\circ a_D= x_D \qquad y\circ b_D= y_D\]
	Let us show that $(x,y)$ is an arrow of $\comma{L}{R}$. Given $D\in \D$ we have
	\begin{align*}
		R(y)\circ f\circ L(a_D)&=R(y)\circ R\left(b_D\right)\circ f_D\\&=R\left(y\circ b_D\right)\circ f_D\\&=R\left(y_D\right)\circ f_D\\&= g\circ L\left(x_D\right)\\&=g\circ L\left(x\circ a_D\right)\\&=g\circ L(x)\circ L\left(a_D\right)
	\end{align*}
	from which it follows  that the following diagram commutes.
	\[\xymatrix@C=35pt{L(A)\ar[d]_{f} \ar[r]^{L(x)} & X \ar[d]^{g}\\ R(B) \ar[r]_{R(y)}& Y}\]
	This shows that $\left((A, B, f), \left\{(a_D, b_D)\right\}_{D\in \D}\right)$ is colimiting for $F$ and the thesis follows.    
\end{proof} 


\Cref{prop:dual} and \Cref{colim} now yields the following.
\begin{corollary}\label[corollary]{lim} The family $\{U_L, U_R\}$ jointly creates limits along every diagram $F\colon \D\to \comma{L}{R}$ such that $R$ preserves the limit of $U_R\circ I$.
\end{corollary}

We can use \Cref{lim} to characterize monos in comma categories. 
\begin{corollary}\label[corollary]{cor:mono}
	If $R$ preserves pullbacks then an arrow $(h,k)$ in $\comma{L}{R}$ is mono if and only if both $h$ and $k$ are monomorphisms.
\end{corollary}
\begin{proof}
	$(\Rightarrow)$  If $(h,k)\colon (A,B,f)\to (A', B', g)$ is a mono then the following square is a pullback in $\comma{L}{R}$
	\[\xymatrix@C=50pt{(A, B, f)  \ar[r]^{\id{(A,B,f)}} \ar[d]_{\id{(A,B,f)}}& (A, B, f)  \ar[d]^{(h,k)} \\ (A, B, f)  \ar[r]_{(h,k)}& (A', B', g) }\]
	Using \Cref{lim} we deduce that the following two squares are pullbacks in $\A$ and $\B$.
	\[\xymatrix{A \ar[r]^{\id{A}} \ar[d]_{\id{A}} & A \ar[d]^{h} & B \ar[r]^{\id{B}} \ar[d]_{\id{B}} & B  \ar[d]^{k}\\ A \ar[r]_{h} & A' & B \ar[r]_{k} & B'}\]
	From which it follows that $h$ and $k$ are monos.
	
	\smallskip\noindent 
	$(\Leftarrow)$ Since $h$ and $k$ are monos then we have two pullback squares 	
	\[\xymatrix{A \ar[r]^{\id{A}} \ar[d]_{\id{A}} & A \ar[d]^{h} & B \ar[r]^{\id{B}} \ar[d]_{\id{B}} & B  \ar[d]^{k}\\ A \ar[r]_{h} & A' & B \ar[r]_{k} & B'}\]
	By \Cref{lim} this implies that 
	\[\xymatrix@C=50pt{(A, B, f)  \ar[r]^{\id{(A,B,f)}} \ar[d]_{\id{(A,B,f)}}& (A, B, f)  \ar[d]^{(h,k)} \\ (A, B, f)  \ar[r]_{(h,k)}& (A', B', g) }\]
	is a pullback in $\comma{L}{R}$ and we are done.   
\end{proof}

We end this section pointing out another useful fact,  showing that in some cases we can guarantee  the existence of a left adjoint to $U_R$. 

\begin{proposition}\label[proposition]{prop:left}
	If $\A$ has initial objects and $L$ preserves them then the forgetful functor $U_R\colon \comma{L}{R}\to \B$ has a left adjoint $\Delta$.
\end{proposition}
\begin{proof} For an object $B\in \B$ we can define $\Delta(B)$ as $(0, B, ?_{B})$, where $0$ is an initial object in $\A$ and $?_{R(B)}$ is the unique arrow $L(0)\to R(B)$. Consider $\id{B}\colon B\to U_R(\Delta(B))$ be the identity, and suppose that a $k\colon B\to U_R(A, B', f)$ in $\B$ is given. By initiality of $0$, there is only one arrow $?_A\colon 0\to A $ in $\A$ and, since $L$ preserves initial objects, the following square commutes.
	\[\xymatrix@C=30pt{L(0) \ar[r]^{L(?_A)} \ar[d]_{?_{R(B)}} & L(A) \ar[d]^{f}\\ R(B) \ar[r]_{R(k)}& R(B')}\]
	Thus $(h,k)$ is the unique morphism $\Delta(B)\to (A, B', f)$	such that $U_R(h,k)=k$.   
\end{proof}
Dualizing we get immediately the following.
\begin{corollary}If $\B$ has terminal objects preserved by $R$ then $U_L\colon \comma{L}{R}\to \A$ has a right adjoint.
\end{corollary}

\subsection{Slice categories}

This section is devoted to recall some basic facts about the so called \emph{slice categories}.
\begin{definition}\index{category!slice -}
	Let $X$ be an object of a category $\X$, we will define the following two categories.
	
	\begin{itemize}
		\item The \emph{ slice category over $X$} is the category $\X/X$ which has as objects arrows $f\colon Y\to X$ and  in which an arrow $h\colon f\to g$ is $h\colon Y\to Y'$ in $\X$ such that the following triangle commutes.
		\[\xymatrix@C=15pt{Y \ar[dr]_{f} \ar[rr]^{h}&& Y' \ar[dl]^{g}\\ & X}\] 
		\item  	 Dually, the \emph{ slice category under $X$} is the category $X/\X$ in which objects are arrows $f\colon X\to Y$  with domain $X$ and a morphism $h\colon  f\to g$ is an arrow of $\X$ fitting in a triangle as the one below.
		\[\xymatrix@C=15pt{&X \ar[dr]^{g}\ar[dl]_{f}\\ Y  \ar[rr]_{h}&& Y'}\] 
	\end{itemize}
\end{definition}

\begin{remark} For every $X\in\ X$ we have forgetful functors
	\[\begin{split}
		\dom_X&\colon \X/X\to \X\\
		\functor[l]{f}{h}{g}
		&\functormapsto
		\rfunctor{\dom(f) }{h}{\dom(g)}
	\end{split}\qquad \begin{split}
		\cod_X&\colon X/\X\to \X\\
		\functor[l]{f}{h}{g}
		&\functormapsto
		\rfunctor{\cod(f) }{h}{\cod(g)}
	\end{split}\]
\end{remark}


We can realize the slice over and under an object $X\in \X$ as comma categories.

\begin{proposition}\label[proposition]{prop:commaapp}
	For every object $X$ in a category $\X$, if $\delta_X\colon \T\to \X$ is  the constant functor of value  $X$ from the category with only one object $*$, then $\X/X$ and $X/\X$ are isomorphic to, respectively,  $\comma{\id{X}}{\delta_X}$ and $\comma{\delta_X}{\id{X}}$ .
\end{proposition}
\begin{proof} Define functors $F_1\colon \comma{\id{X}}{\delta_X}\to \X/X$ and $G_1\colon \X/X\to \comma{\id{X}}{\delta_X}$ as follows
	\[	\begin{split}
		\functor[l]{(Y, *,  f)}{(h, \id{*} )}{( Y', *,  g)}
		\functormapsto
		\rfunctor{f }{h}{g}
	\end{split} \qquad \begin{split}
		\functor[l]{f }{h}{g}
		\functormapsto
		\rfunctor{(\dom(f), *, f)}{(h, \id{*} )}{(\dom(g), *, g)}
	\end{split}\]
	Similarly, we have $F_2\colon \comma{\delta_X}{\id{X}}\to X/\X$ and $G_2\colon X/\X\to \comma{\delta_X}{\id{X}}$
	\[	\begin{split}
		\functor[l]{(*, Y, f)}{(\id{*},h )}{(*, Y', g)}
		\functormapsto
		\rfunctor{f }{h}{g}
	\end{split} \qquad \begin{split}
		\functor[l]{f }{h}{g}
		\functormapsto
		\rfunctor{(*, \cod(f), f)}{(\id{*}, h )}{(*, \cod(g), g)}
	\end{split}\]
	It is now obvious to see that $F_1,G_1$ and $F_2, G_2$ are pairs of inverses.   
\end{proof}

A straightforward application of \Cref{lim} now yields the following.
\begin{corollary}\label[corollary]{cor:pbapp}
	If $\X$ has pullbacks, then for every object $X$, the slice $\X/X$ has pullbacks too.
\end{corollary}




\section{Comparison with Ghica}
\todo{Lo vogliamo fare?}


\end{document}

